\chapimg{fig/sh08.jpg}
\IfFileExists{extra/lead08.tex}{\chapterlead{\subhd{After this chapter you can} classify soil-water models and state the assumptions of the tipping-bucket model; compute runoff by the curve-number method; run a daily bucket by hand and by program; read the dashboard indicators of the application; and design the trigger of a season by comparing simulated seasons.}
}{}
\chapter{The Simulation Engines and the Dashboard}\label{chap:08}
\markboth{The Simulation Engines and the Dashboard}{The Simulation Engines and the Dashboard}
\banner{Overview of the subject}
\IfFileExists{extra/bg08.tex}{\section{Introduction: models of the soil water balance}
Chapter~3 gave the physics of water flow in a soil: the Richards equation. For many purposes (scheduling, design, planning) that is more detail than needed, and a simpler model is used: the soil is a \emph{bucket} whose content changes each day by what comes in and what goes out. Models of the soil water balance can be ranked by their detail. The \emph{tipping bucket} has one store, refilled by rain and irrigation and emptied by evapotranspiration, and spills the excess as drainage. The \emph{cascading bucket} has several stores in series, each spilling to the next. The \emph{capacity models} add a rate of drainage that depends on the content. The \emph{mechanistic} models solve Richards' equation. The application uses two engines: Engine~1, which reproduces the relations of the Daisy ensemble by ratios, and Engine~2, a daily bucket with the FAO stress function; the dashboard presents their results as indicators.

\section{Assumptions and equations of the tipping bucket}
The root zone is one layer of depth $Z_r$ with capacity $TAW$ between field capacity and wilting point. Each day $i$: (1) rain $P_i$ is reduced by interception and runoff to the infiltration $f_i$; (2) irrigation $I_i$ is added; (3) the actual evapotranspiration is $ET_{a,i}=K_{s,i}K_{c,i}ET_{0,i}$; (4) the depletion is updated and the excess over field capacity drains in the same day. The depletion is

\begin{equation}D_i=\min\bigl\{TAW,\ \max\{0,\ D_{i-1}-f_i-I_i+ET_{a,i}\}\bigr\},\qquad DP_i=\max\{0,\ f_i+I_i-ET_{a,i}-D_{i-1}\},\label{eq:bucket8}\end{equation}


with the stress coefficient $K_s=1$ for $D\le RAW$ and $K_s=(TAW-D)/(TAW-RAW)$ beyond it. The model assumes that drainage is instantaneous, that there is no capillary rise and that the root zone is homogeneous; these hold best for deep, freely drained soils.

\section{Runoff: the curve-number method}
Rain that arrives faster than the soil can take it, or on a saturated soil, runs off. The \kw{curve-number} method of the US Soil Conservation Service estimates the direct runoff $Q$ from a storm of depth $P$ by

\begin{equation}Q=\frac{(P-I_a)^2}{P-I_a+S},\quad I_a=0.2S,\quad S=\frac{25\,400}{CN}-254\ \ (\text{mm}),\ \ P>I_a,\label{eq:cn8}\end{equation}


with the curve number $CN$ from 30 (very permeable, dry) to 100 (impervious). A cultivated loam in fair condition has $CN\approx75$ at average antecedent moisture, 87 when the soil is wet and 57 when dry.

\exbox{Runoff of a storm}{A storm of 60 mm falls on a loam with $CN=75$. Find $S$, $I_a$ and the runoff. How does the runoff change if the soil is wet ($CN=87$)?}\sol \givenline{$CN=75$ ; $CN=87$}\ans $S=25400/75-254=84.7$ mm; $I_a=0.2S=16.9$ mm; $Q=(60-16.9)^2/(60-16.9+84.7)=14.5$ mm \ans For $CN=87$: $S=38.0$, $I_a=7.6$, $Q=30.4$ mm \ans Wet soil converts 51 per cent of the storm to runoff, against 24 per cent when it is average. \finans

\section{A season of the bucket, by program}
The bucket is easy to program, and a season can be run in a fraction of a second. The program of this book (and the engine of the application) reads a daily record of $ET_0$ and rain, builds $K_c(t)$ from the four stages, applies the irrigation rule (a trigger on depletion with a dose that refills to field capacity) and writes the totals: $ET_a$, irrigation, drainage, runoff and the relative yield from the Doorenbos--Kassam relation. The result of one season is a number of tables and charts; the design question is what the best trigger is. Section~\ref{sec:proj8} below compares triggers for a maize season.

\section{The dashboard of the application}
The dashboard is the summary page of the application: one set of indicators read from the layers. The most important are (i) the \emph{irrigation} applied, (ii) the \emph{evapotranspiration} and its split into $T$, $E_s$ and $E_o$, (iii) the \emph{drainage} and the nitrogen it carries, (iv) the \emph{yield} and the \emph{water productivity} $WP=Y/ET_a$, (v) the \emph{irrigation efficiency} $IE=(ET_a-ET_{rain})/I$ (the share of the applied water that the crop transpired), and (vi) the \emph{closure} of the balance. Each indicator has a formula, and the chapter on the dashboard shows how every number is produced; here we compute two of them by hand.

\exbox{Dashboard indicators}{A season gives $P=380$, $I=240$, $ET_a=520$, $D=70$, $Q=10$ mm, yield 11.2 Mg\,ha$^{-1}$, and the rainfed twin has $ET_a=410$. Find WP, the irrigation efficiency based on the extra evapotranspiration, and the closure.}\sol \givenline{$P=380$ ; $I=240$ ; $ET_a=520$ ; $D=70$ ; $Q=10$ mm ; $ET_a=410$}\ans $WP=11\,200/5\,200=2.15$ kg\,m$^{-3}$ \ans $IE=(520-410)/240=0.458$ \ans Balance: $P+I=620$; outputs $ET+D+Q=600$, so $\Delta S=20$ mm if the storage has changed by that amount; the residual is zero by definition once $\Delta S$ is entered \ans Only 46 per cent of the irrigation became extra evapotranspiration; the rest drained, ran off or stayed in the soil. \finans

}{}
\banner{The theory step by step}
\IfFileExists{extra/th08.tex}{\section{Meta-models: why a fast model is built from a slow one}
A detailed model such as Daisy takes minutes or hours for a run of twenty years, which is too slow for a decision tool in which the user moves a slider and expects an answer at once. A \kw{meta-model} (surrogate, emulator) is a simple model that is fitted to the results of the detailed one, and then replaces it inside the tool. Two ways are used. The first is a \emph{statistical emulator}, a regression or a neural network of the outputs on the inputs; it is fast but gives no insight and fails outside the range of the training data. The second, used here, is a \emph{physically structured meta-model}: a water balance with few parameters in which the results of Daisy supply the ratios and the scaling factors. It keeps the conservation of water by construction, and the parameters have a meaning.
\par\noindent \textbf{Note:} A meta-model must be tested on runs that were not used to fit it (held-out runs), and it must say when the user is outside the range for which it was fitted.

\section{The bucket model of the soil water}
The simplest model of the soil water treats the root zone as a \kw{bucket}: it is filled by rain and irrigation, emptied by evapotranspiration, and overflows as drainage when it is full. In daily steps the storage $S$ is
\begin{equation*}\boxed{S_{d+1}=\min\bigl(S_{max},\ S_d+P_d+I_d-ET_{a,d}\bigr),\qquad D_d=\max\bigl(0,\ S_d+P_d+I_d-ET_{a,d}-S_{max}\bigr)}\end{equation*}
where $S_{max}$ is the capacity of the bucket (field capacity), $P$ the rain, $I$ the irrigation, $ET_a$ the actual evapotranspiration and $D$ the drainage. With a second bucket for the layer below it is called a \kw{tipping-bucket} (cascade) model. The model has two advantages (it conserves water exactly and needs few parameters) and one limit: it has no suction, so that it cannot give the trigger directly, and the drainage is instant, so it has no memory of the rate of flow. The engines of this chapter add the suction from the retention curve and the delay from the conductivity.

\section{Runoff: the curve-number method}
Where rain falls faster than the soil can take it, water runs off. The method of the Soil Conservation Service of the USA estimates the runoff $Q$ from the rain $P$ (mm) of an event by
\begin{equation*}\boxed{Q=\frac{(P-0.2S)^2}{P+0.8S}\quad(P>0.2S),\qquad S=\frac{25\,400}{CN}-254}\end{equation*}
where $S$ is the potential retention (mm) and $CN$ the curve number, from 30 (permeable, covered soil) to 100 (impervious surface). It is empirical and has no time scale; for a CN of 75 and an event of 50 mm, $S=84.7$ mm and $Q=(50-16.9)^2/(50+67.7)=9.3$ mm. The explicit solution of the infiltration (Green--Ampt, Chapter~\ref{chap:11}) is used where the intensity matters.

\section{Stochastic weather: how rain is generated}
When the weather of a future or a different climate is needed, a \kw{weather generator} makes a series of days with the right statistics. Rain is made in two steps. The \emph{occurrence} is a first-order Markov chain: the probability that a day is wet depends on whether the previous day was wet or dry, with two parameters $p_{WW}$ and $p_{DW}$. The \emph{amount} of a wet day follows a skewed distribution, usually the gamma distribution with a shape and a scale parameter, fitted to the record. Temperature and radiation follow with their seasonal means and their dependence on whether the day is wet. The generated series reproduces the mean, the variance and the frequency of wet spells of the record, but not the sequence of the years.

\section{Closure of the budget as a test}
A water-balance model must close: the change of storage equals the inputs less the outputs, to machine accuracy. The \kw{closure residual} is a test of the program, not of the physics. A residual of $10^{-12}$ mm means that no water is created or destroyed in the code; a residual that grows with time points to a bug in one of the fluxes. Section~\ref{sec:closure} shows the closure of the engines.
}{}
\banner{Engine 1: depletion, critical suction and the meta-model}\label{ch:e1}
Part II treats the application layer by layer. Every layer of SWAFL is driven by one of two engines (the Daisy meta-model and the daily bucket model) or by a design module. For each layer this part states the relations that the application gives, then solves them with numbers. The numbers come from running the stated equations; every problem is solved in full, and where an input is not published (an annual flux, say) it is labelled as an assumed value.

\section{Depletion between field capacity and wilting point}
The application defines depletion of plant-available water between field capacity (suction 100 cm, pF 2.0) and the wilting point (15,849 cm, pF 4.2):
\begin{equation}\label{eq:fdep}\boxed{f(h)=\frac{\Se(100)-\Se(|h|)}{\Se(100)-\Se(15\,849)},\qquad f_{\rm used}=\min\{1,\max\{0,f\}\}.}\end{equation}
\par\noindent where $f$ is the depletion fraction (0 at field capacity, 1 at wilting point), $S_e(h)$ the effective saturation at suction $h$, $|h|=100$ cm the field-capacity suction, $|h|=15\,849$ cm the wilting-point suction and $f_{\rm used}$ the value limited to $[0,1]$.

$f=0$ is field capacity and $f=1$ is the wilting point.

\exbox{Depletion at 300 cm in three soils}{Compute $f$ at the sensor suction $|h|=300$ cm for the loam S06 ($\alpha=0.036$, $n=1.56$), the sand S03 ($\alpha=0.145$, $n=2.68$) and the silty clay loam S10 ($\alpha=0.010$, $n=1.23$).}\sol \givenline{$|h|=300$ cm ; $\alpha=0.036$ ; $n=1.56$ ; $\alpha=0.145$ ; $n=2.68$ ; $\alpha=0.010$ ; $n=1.23$}\ans For each soil compute $m=1-1/n$ and three values of $\Se=[1+(\alpha h)^n]^{-m}$. \finans

\subhd{Loam} ($m=0.3590$): $\Se(100)=0.4663$, $\Se(300)=0.2615$, $\Se(15849)=0.0286$. Then $f=(0.4663-0.2615)/(0.4663-0.0286)=0.468$.

\subhd{Sand} ($m=0.6269$): $\Se(100)=0.0112$, $\Se(300)=0.0018$, $\Se(15849)=0.0000$. Then $f=(0.0112-0.0018)/(0.0112-0.0000)=0.842$.

\subhd{Silty clay loam} ($m=0.1870$): $\Se(100)=0.8784$, $\Se(300)=0.7440$, $\Se(15849)=0.3118$. Then $f=(0.8784-0.7440)/(0.8784-0.3118)=0.237$.

\subhd{Reading.} At $-300$ cm the sand has used 84 per cent of its plant-available water while the silty clay loam has used 24 per cent. The same suction means very different dryness in different soils; this is the reason for the trigger law of Chapter~\ref{ch:ctrl}.


\exbox{Inverting the depletion by bisection}{Find the suction at which the loam S06 reaches the soil-independent critical depletion $f^*=0.45$. The application solves $f(h)=f^*$ by bisection on $\log_{10}|h|\in[0,4.6]$ with 60 halvings.}\sol \givenline{$f^*=0.45$}\ans Start with the interval $[0,4.6]$. At each step take the midpoint $x=\log_{10}h$, evaluate $f$ with Equation~(\ref{eq:fdep}) and keep the half that still brackets $0.45$ (if $f<0.45$ the root is above, so the lower end moves up). The first eight steps are in Table~\ref{tab:bis}; the interval has then shrunk to $4.6/2^8=0.018$ in $\log_{10}h$, that is 4 per cent in $h$. After 60 halvings the width is $4\times10^{-18}$ and the answer is $h^*=284.3$ cm.
\begin{table}[!htbp]\centering\small\caption{First eight bisection steps for the loam at $f^*=0.45$.}\label{tab:bis}
\begin{tabular}{lrrrr}\toprule
\hd{Step} & \hd{$\log_{10}h$ mid} & \hd{$h$ (cm)} & \hd{$f(h)$} & \hd{Keep}\\\midrule
1 & 2.3000 & 199.5 & 0.3201 & lower\\
2 & 3.4500 & 2818.4 & 0.8935 & upper\\
3 & 2.8750 & 749.9 & 0.7053 & upper\\
4 & 2.5875 & 386.8 & 0.5456 & upper\\
5 & 2.4437 & 277.8 & 0.4422 & lower\\
6 & 2.5156 & 327.8 & 0.4961 & upper\\
7 & 2.4797 & 301.8 & 0.4697 & upper\\
8 & 2.4617 & 289.5 & 0.4561 & upper\\\bottomrule\end{tabular}\end{table}
\textbf{Check.} The application's table gives 284 cm for the loam; the computed value agrees to the rounding of the printed parameters. \finans


\begin{kfig}{Depletion $f$ of plant-available water against suction: zero at field capacity (100 cm), one at the wilting point (15\,849 cm); the loam S06 is drawn.}
\draw[axis] (0,0)--(7.4,0);\draw[axis] (0,0)--(0,4.4);
\node[lab,below] at (3.7,-0.55) {Suction $|h|$ (cm, log scale)};\node[lab,rotate=90] at (-0.8,2.2) {Depletion $f$};
\foreach \x/\t in {1.5/10,3/100,4.5/1000,6/10000}{\draw[ln] (\x,0)--(\x,-0.1);\node[lab,below] at (\x,-0.1) {\t};}
\foreach \y/\t in {0/0,2/0.5,4/1}{\draw[ln] (0,\y)--(-0.1,\y);\node[lab,left] at (-0.1,\y) {\t};}
\draw[crv,smooth] plot coordinates {(3.000,0.000) (3.055,0.175) (3.110,0.345) (3.165,0.510) (3.220,0.670) (3.275,0.825) (3.330,0.974) (3.385,1.118) (3.440,1.256) (3.495,1.389) (3.550,1.517) (3.605,1.639) (3.660,1.757) (3.715,1.870) (3.770,1.978) (3.825,2.081) (3.880,2.180) (3.935,2.275) (3.990,2.366) (4.045,2.452) (4.100,2.535) (4.155,2.614) (4.210,2.689) (4.265,2.761) (4.320,2.830) (4.375,2.896) (4.430,2.959) (4.485,3.019) (4.540,3.076) (4.595,3.130) (4.650,3.183) (4.705,3.232) (4.760,3.280) (4.815,3.325) (4.870,3.368) (4.925,3.409) (4.980,3.449) (5.035,3.486) (5.090,3.522) (5.145,3.556) (5.200,3.589) (5.255,3.620) (5.310,3.649) (5.365,3.677) (5.420,3.704) (5.475,3.730) (5.530,3.755) (5.585,3.778) (5.640,3.800) (5.695,3.822) (5.750,3.842) (5.805,3.861) (5.860,3.880) (5.915,3.897) (5.970,3.914) (6.025,3.930) (6.080,3.946) (6.135,3.960) (6.190,3.974) (6.245,3.987) (6.300,4.000)};
\draw[gd] (3,0)--(3,4.2);\draw[gd] (6.3,0)--(6.3,4.2);\node[lab,above] at (3,4.2) {FC};\node[lab,above] at (6.3,4.2) {WP};
\end{kfig}

\section{Table of critical suctions for the ten soils}
\exbox{The table of critical suctions}{Reproduce the application's table of effective parameters: for each of the ten soils compute the suction $h^*$ at $f^*=0.45$ and the depletion at 300 cm.}\sol \givenline{$f^*=0.45$}\ans Repeat the bisection for every soil. The result is in Table~\ref{tab:h45}. The published values are $387, 262, 143, 160, 195, 284, 401, 311, 503, 729$ cm and $0.380, 0.496, 0.842, 0.756, 0.627, 0.468, 0.367, 0.439, 0.314, 0.237$; the largest difference is 1 cm and 0.001 in $f$, which comes from the rounding of the printed parameters.
\begin{table}[!htbp]\centering\small\caption{Critical suction at $f^*=0.45$ and depletion at 300 cm, computed from Equation~(\ref{eq:fdep}).}\label{tab:h45}
\begin{tabular}{lrrrrrr}\toprule
\hd{ID} & \hd{Soil} & \hd{$\alpha$ (cm$^{-1}$)} & \hd{$n$} & \hd{AWC (mm/cm)} & \hd{$h^*$ (cm)} & \hd{$f(300)$}\\\midrule
S01 & Coarse loam & 0.0961 & 1.353 & 1.5 & 387 & 0.379\\
S02 & Coarse sand & 0.0536 & 1.600 & 0.5 & 262 & 0.496\\
S03 & Sand & 0.1450 & 2.680 & 0.6 & 143 & 0.842\\
S04 & Loamy sand & 0.1240 & 2.280 & 0.9 & 160 & 0.756\\
S05 & Sandy loam & 0.0750 & 1.890 & 1.2 & 195 & 0.627\\
S06 & Loam & 0.0360 & 1.560 & 1.7 & 284 & 0.468\\
S07 & Silt loam & 0.0200 & 1.410 & 2.0 & 401 & 0.366\\
S08 & Sandy clay loam & 0.0590 & 1.480 & 1.4 & 311 & 0.438\\
S09 & Clay loam & 0.0190 & 1.310 & 1.6 & 503 & 0.314\\
S10 & Silty clay loam & 0.0100 & 1.230 & 1.7 & 729 & 0.237\\\bottomrule\end{tabular}\end{table}
\textbf{Reading.} The critical suction rises from 143 cm for the sand to 729 cm for the silty clay loam: a fivefold spread in suction for one and the same depletion. A trigger at a fixed suction is therefore a different management in every soil. \finans


\exbox{Net irrigation depth and rooting depth}{The default net irrigation depth is $Z_n=\mathrm{clip}(f_T\cdot\mathrm{AWC}\cdot R_{\max};25,150)$ mm, with the maximum rooting depth
\begin{equation}\label{eq:rmax}\boxed{R_{\max}=\mathrm{clip}\Bigl(R_c(0.85+0.3\,\mathrm{sand})\bigl[1+0.12\log_{10}(|h_T|/300)\bigr]\bigl[1-0.18\tfrac{N-100}{100}\bigr];25,100\Bigr)\ \text{cm}.}\end{equation}
\par\noindent where $R_{\max}$ is the maximum rooting depth (cm), $R_c$ the value for the crop class (cm), ${\rm sand}$ the sand fraction, $h_T$ the trigger suction (cm) and $N$ the nitrogen supply (kg\,ha$^{-1}$).

For a cereal $R_c=95$ cm. Find $Z_n$ for the loam S06 (sand 40 per cent, AWC 1.7 mm cm$^{-1}$) with trigger $-300$ cm and nitrogen $100$ per cent.}\sol \givenline{$Z_n=\mathrm{clip}(f_T\cdot\mathrm{AWC}\cdot R_{\max};25,150)$ mm ; $R_c=95$ cm}\ans 1. Rooting depth: $0.85+0.3(0.40)=0.97$; the suction factor is $1+0.12\log_{10}(300/300)=1$; the nitrogen factor is $1-0.18(0)=1$. So $R_{\max}=95(0.97)=92.15$ cm, inside the clip $[25,100]$. \finans

2. Trigger depletion: $f_T=f(300)=0.468$ (Table~\ref{tab:h45}).

3. Depth: $Z_n=0.468\times1.7\times92.15=73.3$ mm, inside $[25,150]$ mm.

\subhd{Repeat for the sand S03} (sand 92 per cent). $R_{\max}=95(0.85+0.276)=106.97$, clipped to $100$ cm. $f_T=0.842$, so $Z_n=0.842\times0.6\times100=50.5$ mm. The sand needs a smaller net depth per event (50.5 mm against 73.3 mm) because it holds less water per centimetre, although it is irrigated from a deeper depletion.

\subhd{Effect of a drier trigger.} For the loam at $h_T=1000$ cm: the suction factor is $1+0.12\log_{10}(3.333)=1.0627$, $R_{\max}=92.15\times1.0627=97.9$ cm, $f_T=f(1000)$ is $0.759$, so $Z_n=0.759\times1.7\times97.9=126.3$ mm. A drier trigger deepens the roots by 6 per cent and increases the depth of each application by 72 per cent.


\section{The meta-model: ratios of Daisy results}
The annual fluxes of Engine 1 are built from the 20-year Daisy means $\mathbf b=(E_b,T_b,ET_b,I_b,Y_b,D_b,N_b)$ and from lookups in the mulch tables $M(\varphi,\rho)$ and warming tables $W(\Delta T,\rho)$ of the ensemble. The structure is multiplicative for fluxes and additive for drainage:
\begin{align}
T&=T_b\Bigl[\frac{\mathrm{now}.T}{\mathrm{bare}.T}\Bigr]\,\mathrm{clip}\Bigl(\frac{\mathrm{tN}.T}{\mathrm{tB}.T};0.3,2.5\Bigr),\label{eq:mT}\\
D&=D_b+(\mathrm{now}.D-\mathrm{bare}.D)+(\mathrm{tN}.D-\mathrm{tB}.D).\label{eq:mD}
\end{align}
Fluxes are multiplied because ratios of large numbers are stable; drainage and leaching are added because ratios of numbers near zero are not.

\exbox{Combining mulch and warming ratios}{For a soil the baseline is $T_b=240$ mm, $D_b=40$ mm. The mulch table gives, at the present mulch and rain, $\mathrm{now}.T=232$, $\mathrm{bare}.T=225$, $\mathrm{now}.D=55$, $\mathrm{bare}.D=38$; the warming table gives $\mathrm{tN}.T=236$, $\mathrm{tB}.T=229$, $\mathrm{tN}.D=44$, $\mathrm{tB}.D=40$. (These table entries are \emph{assumed} illustration values; the real ones are in the ensemble tables.) Find $T$ and $D$.}\sol \givenline{$T_b=240$ mm ; $D_b=40$ mm ; $\mathrm{now}.T=232$ ; $\mathrm{bare}.T=225$ ; $\mathrm{now}.D=55$ ; $\mathrm{bare}.D=38$ ; ; $\mathrm{tN}.T=236$ ; $\mathrm{tB}.T=229$ ; $\mathrm{tN}.D=44$}\ans Transpiration: the mulch ratio is $232/225=1.0311$; the warming ratio is $236/229=1.0306$, inside $[0.3,2.5]$. So $T=240\times1.0311\times1.0306=255.0$ mm. \finans

Drainage: $D=40+(55-38)+(44-40)=40+17+4=61$ mm.

\subhd{Why not multiply drainage?} If $\mathrm{bare}.D$ were $2$ mm and $\mathrm{now}.D$ were $6$ mm, the ratio would be 3 and $D_b=40$ would become 120 mm; the additive form adds only the 4 mm difference. This is why the application treats drainage and leaching as differences.


\exbox{Closure window and irrigation efficiency}{After the four corrections the model enforces closure: $ET\in[P+I-D-30,\ P+I-D+15]$; if $ET$ lies outside, $ET$ is clipped and $E$, $T$ are scaled by the same factor. Then $E_{\rm other}=\max(0,ET-T-E)$, $\Delta S=P+I-ET-D$ and the irrigation efficiency $\mathrm{IE}=100Y/I$ (kg m$^{-3}$). Take $P=600$, $I=150$, $D=61$, $T=255$, $E=190$, $ET=520$ and $Y=6.4$ Mg ha$^{-1}$.}\sol \givenline{$E_{\rm other}=\max(0,ET-T-E)$ ; $\mathrm{IE}=100Y/I$ ; $P=600$ ; $I=150$ ; $D=61$ ; $T=255$ ; $E=190$ ; $ET=520$ ; $Y=6.4$ Mg ha$^{-1}$}\ans 1. Window: $P+I-D=689$, so $ET$ must lie in $[659,704]$. The model value $ET=520$ is below the lower bound. \finans

2. Clip: $ET\leftarrow659$. The scale factor is $659/520=1.2673$, so $T=255\times1.2673=323.2$ and $E=190\times1.2673=240.8$ mm.

3. $E_{\rm other}=\max(0,659-323.2-240.8)=95.0$ mm (interception and mulch evaporation).

4. $\Delta S=600+150-659-61=30$ mm.

5. $\mathrm{IE}=100(6.4)/150=4.27$ kg m$^{-3}$ (1 mm is 10 m$^3$ ha$^{-1}$, so 150 mm is 1,500 m$^3$ and 6,400 kg / 1,500 m$^3$ = 4.27).

\subhd{Reading.} The window protects the balance from independent effects compounding beyond what the water allows; here it raised $ET$ by 139 mm, a large correction, which would signal that the ratio inputs are inconsistent.


\section{Sensitivity of the reported numbers}
\exbox{Sensitivity of the depth to the water capacity}{For the loam, $Z_n$ is proportional to $f_T$, AWC and $R_{\max}$. Compute the relative change of $Z_n$ if AWC is overestimated by 20 per cent.}\sol  $Z_n\propto \mathrm{AWC}$, so the relative change is $+20$ per cent: $73.3\to88.0$ mm. Because $f_T=f(h_T)$ does not involve AWC, the suction at which irrigation starts is unchanged; only the depth per application changes. The number of events per season falls in proportion, so the seasonal total is first-order unaffected, but the events are larger and rarer.
\banner{Engine 2: the daily bucket model, step by step}\label{ch:e2}
The daily model has $N=40$ layers of thickness $\Delta z=2.5$ cm. Layer $i$ holds water $w_i$ between $0$ and $W_{\max}=\mathrm{AWC}\,\Delta z$ (mm); relative content $s_i=w_i/W_{\max}$ and depletion $f_i=1-s_i$. All layers start at $0.72W_{\max}$. The sensor layer is $j_s=\mathrm{round}(z_s/\Delta z)$. \finans

\exbox{Layout and state}{For the loam (AWC $=1.7$ mm cm$^{-1}$) with sensor depth 20 cm find $W_{\max}$, the sensor layer, the starting water and the trigger depletion for $-300$ cm.}\sol \givenline{$=1.7$ mm cm$^{-1}$}\ans $W_{\max}=1.7\times2.5=4.25$ mm; $j_s=\mathrm{round}(20/2.5)=8$; starting water $0.72\times4.25=3.06$ mm per layer, so the 100 cm profile holds $40\times3.06=122.4$ mm of plant-available water out of $170$ mm. The trigger depletion is $f_T=0.468$ (Chapter~\ref{ch:e1}), so the sensor layer must fall to $s_8\le1-0.468=0.532$, that is $w_8\le2.26$ mm, to start irrigation. At the start $w_8=3.06>2.26$, so no irrigation is due. \finans


\section{Crop calendar and canopy cover}
For spring barley the sowing day is 100 and the harvest day is 225. The cover is the logistic
\begin{equation}\boxed{c(d)=\frac{1}{1+\exp[-(d-d_{\rm mid})/13]},\quad d_{\rm mid}=d_{\rm sow}+50,}\end{equation}
\par\noindent where $c(d)$ is the canopy cover (0 to 1) on day $d$, $d_{\rm mid}$ the day of half cover, $d_{\rm sow}$ the day of sowing and 13 the width of the logistic rise (d).

which after the peak day $d_{\rm pk}=d_{\rm harv}-50$ declines linearly to zero at harvest: $c\leftarrow c\,\max\{0,1-(d-d_{\rm pk})/(d_{\rm harv}-d_{\rm pk})\}$.

\exbox{Canopy cover through the season}{Compute the canopy cover of spring barley on days 100, 120, 140, 150, 160, 175, 190, 205, 215 and 225.}\sol  $d_{\rm mid}=150$, $d_{\rm pk}=175$. Before the peak $c=1/(1+e^{-(d-150)/13})$: on day 120, $(120-150)/13=-2.308$, $e^{2.308}=10.05$, so $c=1/11.05=0.0905$; on day 150, $c=0.5$; on day 160, $e^{-0.769}=0.4634$ so $c=0.683$. At the peak $d=175$: $e^{-1.923}=0.1462$, $c=0.872$. After the peak the factor is $1-(d-175)/50$: on day 190 the logistic gives $0.956$ and the factor is $1-15/50=0.70$, so $c=0.669$; on day 205, $0.986\times0.40=0.394$; on day 225, the factor is $0$ and $c=0$. Table~\ref{tab:can} lists all ten days.\begin{table}[!htbp]\centering\small\caption{Canopy cover of spring barley.}\label{tab:can}
\begin{tabular}{lr}\toprule
\hd{Day} & \hd{Cover $c$}\\\midrule
100 & 0.021\\
120 & 0.090\\
140 & 0.317\\
150 & 0.500\\
160 & 0.683\\
175 & 0.872\\
190 & 0.669\\
205 & 0.394\\
215 & 0.199\\
225 & 0.000\\\bottomrule\end{tabular}\end{table}
\textbf{Reading.} The curve is engineered to reproduce the shape of a canopy, not a measured leaf-area record; the cover mean over the year, $\bar C=\frac1{365}\sum c(d)=0.138$ is the divisor that spreads the annual transpiration over the days. \finans


\section{Rainfall generator}
Rain falls on a day with probability 0.28; the wet-day depth is exponential and capped,
\begin{equation}\boxed{R=\min\{70,\ \mu_0\gamma_y(-\ln U)\},\ U\sim\mathrm{Uniform}(0,1),\quad \mu_0=\frac{P}{0.28\times365}.}\end{equation}
\par\noindent where $R$ is the daily rainfall (mm), $U$ a random number uniform on (0,1), $\mu_0$ the mean depth of a rainy day (mm), $\gamma_y$ the scale of year $y$, $P$ the annual rainfall (mm) and 0.28 the assumed probability of a rainy day.

The canopy and litter catch the first $I_c=\mu_0\ln(1/r)$ mm of every shower, with $r=\max\{40,P-P_{\rm other}\}/P$, and only throughfall $T_f=\max\{0,R-I_c\}>0.3$ mm enters the soil.

\exbox{Rainfall generator constants}{For an annual rain $P=600$ mm, find $\mu_0$, the expected number of wet days, the probability that a wet day exceeds the cap, and the interception $I_c$ if $P_{\rm other}=60$ mm.}\sol \givenline{$P=600$ mm ; $P_{\rm other}=60$ mm}\ans $\mu_0=600/(0.28\times365)=5.871$ mm. Expected wet days: $0.28\times365=102.2$. The exponential tail gives $P(R>70)=\exp(-70/5.871)=\exp(-11.92)=6.7\times10^{-6}$ (with $\gamma_y=1$), so the cap hardly binds in a year with 102 wet days (the chance that it binds at least once is $102\times6.7\times10^{-6}=7\times10^{-4}$). Next, $r=\max\{40,600-60\}/600=540/600=0.9$ and $I_c=5.871\ln(1/0.9)=0.619$ mm. \finans

\subhd{Check.} The expected interception per wet day is $\mu_0(1-r)=5.871\times0.1=0.587$ mm, so the yearly total is $102.2\times0.587=60.0$ mm, which equals $P_{\rm other}$. The choice of $I_c$ makes the daily generator reproduce the annual interception of Engine 1.


\exbox{The year factor}{The dryness factor is $\gamma_y=\mathrm{clip}\bigl((m_{\rm med}/m_y)^{0.3};0.65,1.45\bigr)$ with $m_y=10^{\overline{\log_{10}h}}$ the geometric-mean suction of the growing season. If the median season has $m_{\rm med}=300$ cm and year $y$ has $m_y=900$ cm, find $\gamma_y$ and the mean wet-day depth of that year; do the same for $m_y=60$ cm.}\sol \givenline{$m_y=10^{\overline{\log_{10}h}}$ ; $m_{\rm med}=300$ cm ; $m_y=900$ cm ; $m_y=60$ cm}\ans Dry year: $(300/900)^{0.3}=0.3333^{0.3}=e^{0.3\ln0.3333}=e^{-0.3296}=0.719$; the clip does not bind, so $\gamma_y=0.719$ and the mean wet-day depth is $5.871\times0.719=4.22$ mm. Wet year: $(300/60)^{0.3}=5^{0.3}=e^{0.4828}=1.621$, clipped to $1.45$, so the mean depth is $5.871\times1.45=8.51$ mm. A drier record (higher suction) gives drier generated years, and the clip stops the factor from exceeding $0.65$ to $1.45$. \finans


\section{Infiltration, redistribution and drainage}
Water enters the top layer and cascades: any layer above $W_{\max}$ passes the excess down; the excess leaving layer 40 is drainage. Between layers a slow redistribution moves water toward the drier neighbour, $q_i=0.08(w_{i+1}-w_i)$, $w_i\leftarrow w_i+q_i$, $w_{i+1}\leftarrow w_{i+1}-q_i$.

\exbox{A 10 mm shower in the loam}{The loam ($W_{\max}=4.25$ mm) has all layers at $3.06$ mm. A throughfall of $10.0$ mm arrives. Find the water in layers 1 to 4 after the cascade, then after one redistribution pass over the pair (1,2).}\sol \givenline{$W_{\max}=4.25$ mm}\ans Layer 1: $3.06+10.00=13.06$ mm, an excess of $13.06-4.25=8.81$ mm, so $w_1=4.25$ and $8.81$ mm moves down. Layer 2: $3.06+8.81=11.87$, excess $7.62$, $w_2=4.25$. Layer 3: $3.06+7.62=10.68$, excess $6.43$, $w_3=4.25$. Layer 4: $3.06+6.43=9.49$, excess $5.24$, $w_4=4.25$. The front continues until the excess is used up: each layer absorbs $4.25-3.06=1.19$ mm, so $10.00/1.19=8.4$ layers are filled, that is the wetting front reaches $8.4\times2.5=21$ cm. No drainage occurs at 100 cm. \finans

\subhd{Redistribution.} Within the filled layers $w_i=w_{i+1}$ so $q=0$. At the front, layers 8 and 9 have $w_8=4.25$ and $w_9=3.06+0.4\times1.19=3.54$: $q_8=0.08(3.54-4.25)=-0.057$ mm, so $0.057$ mm moves from layer 8 to layer 9 (downward), smoothing the front by 0.057 mm per day per layer pair.


\section{Roots}
The rooting depth chases water while the canopy grows,
\begin{equation}R_d\leftarrow\max\Bigl\{R_d,\ \min\Bigl\{R_{\max}\,\mathrm{clip}\bigl(\tfrac{c-0.04}{0.96};0,1\bigr)^{0.8},\ R_{\max},\ R_d+1.7(0.35+\bar s_{\rm wet})\kappa\Bigr\}\Bigr\},\end{equation}
\par\noindent where $R_d$ is the rooting depth on day $d$ (cm), $R_{\max}$ the maximum depth, $c$ the canopy cover, $\bar s_{\rm wet}$ the mean wetness of the soil profile (0 to 1) and $\kappa$ the soil resistance factor; the depth never decreases and grows by at most the third term.

and the root density is
\begin{equation}\boxed{\rho_i\propto\exp\Bigl(-\frac{z_i}{\lambda R_d+8}\Bigr)(0.45+0.9s_i)\Bigl[1+3e^{-((z_i-z_d)/12)^2}\Bigr]_{\rm drip\ only},\quad z_i\le R_d,\ \sum\rho_i=1,}\end{equation}
with $\lambda=0.32$ for cereal.

\exbox{Rooting depth advance}{On day 150 ($c=0.5$) the root depth is $R_d=20$ cm, $R_{\max}=92.15$ cm, the mean relative content under the front is $\bar s_{\rm wet}=0.60$ and the topsoil is wet ($\kappa=1$). Find the new $R_d$.}\sol \givenline{$c=0.5$ ; $R_d=20$ cm ; $R_{\max}=92.15$ cm ; $\bar s_{\rm wet}=0.60$ ; $\kappa=1$}\ans Three candidates: (a) the cover limit $92.15\,\mathrm{clip}((0.5-0.04)/0.96)^{0.8}=92.15\times0.4792^{0.8}=92.15\times0.5549=51.2$ cm; (b) $R_{\max}=92.15$ cm; (c) the growth rate $20+1.7(0.35+0.60)(1)=20+1.615=21.61$ cm. The minimum is 21.61 cm, and the maximum with the old value $20$ is 21.61 cm. The root front therefore advances by 1.6 cm per day while the canopy supports a deeper root: at this rate it needs about $(92-20)/1.6=45$ days to reach full depth. \finans


\exbox{Root density by layer}{With $R_d=15$ cm and the six layers of 0 to 15 cm having relative contents $s=0.50,0.55,0.60,0.65,0.70,0.70$ (no drip), compute the root densities.}\sol \givenline{$R_d=15$ cm ; $s=0.50,0.55,0.60,0.65,0.70,0.70$}\ans $\lambda R_d+8=0.32(15)+8=12.8$ cm. Layer centres $z_i=1.25,3.75,\dots,13.75$ cm. The product $\exp(-z/12.8)(0.45+0.9s)$ is evaluated in Table~\ref{tab:rho} and normalised to sum to 1 (the total before normalisation is $3.4686$). Root density decreases with depth by the exponential and increases with wetness by the second factor; here the wetter lower layers compensate only part of the exponential decline.\begin{table}[!htbp]\centering\small\caption{Root density of the first six layers.}\label{tab:rho}
\begin{tabular}{lrrrr}\toprule
\hd{$z_i$ (cm)} & \hd{$s_i$} & \hd{$e^{-z/12.8}$} & \hd{$0.45+0.9s$} & \hd{$\rho_i$}\\\midrule
1.25 & 0.50 & 0.9070 & 0.900 & 0.2353\\
3.75 & 0.55 & 0.7460 & 0.945 & 0.2033\\
6.25 & 0.60 & 0.6137 & 0.990 & 0.1752\\
8.75 & 0.65 & 0.5048 & 1.035 & 0.1506\\
11.25 & 0.70 & 0.4152 & 1.080 & 0.1293\\
13.75 & 0.70 & 0.3416 & 1.080 & 0.1064\\\bottomrule\end{tabular}\end{table} \finans


\section{Evaporation and transpiration}
Daily soil evaporation spreads the annual $E_{\rm yr}$ with a sinusoidal weight reduced by canopy shade,
\begin{equation}\boxed{E_d=E_{\rm yr}\frac{w({\rm doy})}{\sum_{k=0}^{364}w(k)}(1-0.3c),\quad w({\rm doy})=0.35+0.65\cdot\tfrac12\Bigl[1+\sin\tfrac{2\pi({\rm doy}-100)}{364}\Bigr],}\end{equation}
\par\noindent where $E_d$ is the daily evaporation (mm), $E_{\rm yr}$ the annual total (mm), $w({\rm doy})$ the seasonal weight on the day of the year, $c$ the canopy cover (0 to 1), which shades the soil, and the sum in the denominator normalises the weights.

drawn from the top four layers with weights $(0.55,0.25,0.12,0.08)$, never below $0.05W_{\max}$. Transpiration demand is $T_d=k_tT_{\rm yr}c(d)/(\bar C\cdot365)$ divided among layers by $q_i^T=T_d\rho_ia_i/\sum_j\rho_ja_j$ with $a_i=\mathrm{clip}\bigl((s_i-0.15)/0.6;0,1\bigr)$.

\exbox{Evaporation of one day}{With $E_{\rm yr}=170$ mm, find the evaporation on day 181 (doy index 180), when the canopy cover is $c=0.8$, and how it is drawn from the layers.}\sol \givenline{$E_{\rm yr}=170$ mm ; $c=0.8$}\ans The seasonal weight is $w(180)=0.35+0.65\times0.5[1+\sin(2\pi\,80/364)]=0.35+0.325[1+0.9820]=0.9942$. The sum of the weights over the year is $\sum w=246.05$ (mean 0.674). Hence $E_d=170\times0.9942/246.05\times(1-0.24)=0.522$ mm. It is taken from layers 1 to 4 as $0.55,0.25,0.12,0.08$ times this: $0.287,0.131,0.063,0.042$ mm, provided that each layer stays above $0.05W_{\max}=0.21$ mm. Without a canopy ($c=0$) the same day would lose $0.687$ mm, 30 per cent more than 0.522: shade reduces evaporation. \finans


\exbox{Transpiration demand and its division}{With $T_{\rm yr}=230$ mm, $\bar C=0.138$ and $k_t=1$, find $T_d$ on a day with $c=0.8$ and distribute it among three layers with $\rho=(0.5,0.3,0.2)$ and $s=(0.8,0.5,0.2)$.}\sol \givenline{$T_{\rm yr}=230$ mm ; $\bar C=0.138$ ; $k_t=1$ ; $c=0.8$ ; $\rho=(0.5,0.3,0.2)$ ; $s=(0.8,0.5,0.2)$}\ans $T_d=230\times0.8/(0.138\times365)=184/50.37=3.653$ mm. Availability $a_i=\mathrm{clip}((s_i-0.15)/0.6)$: $(0.8-0.15)/0.6=1.083\to1$; $(0.5-0.15)/0.6=0.583$; $(0.2-0.15)/0.6=0.083$. The weights $\rho_ia_i=0.5,\ 0.175,\ 0.0167$ sum to $0.6917$. The shares are $0.723,\ 0.253,\ 0.024$, so the uptake is $2.641,\ 0.924,\ 0.088$ mm. The wettest layer supplies 72 per cent although it holds only half of the roots, and the dry third layer is almost shut off: this is the linear stress function between 15 and 75 per cent of plant-available water. \finans


\section{The irrigation rule and a season run}
An irrigation of dose $d$ is applied on a day if the canopy cover exceeds 0.2, $95\le\mathrm{doy}\le250$, the spacing and budget rules hold, the sensor layer is dry enough ($1-s_j\ge f_T$) and $f_T<0.995$. The dose is $d=\mathrm{clip}(\bar I/10;15,35)$ mm.

\exbox{One season of the loam}{For $\bar I=150$ mm yr$^{-1}$ find the dose; then run the daily model for one season of spring barley on the loam with $P=600$ mm, $E_{\rm yr}=170$, $T_{\rm yr}=230$, trigger $-300$ cm, sensor at 20 cm, $k_t=1$ and a fixed random seed. (Simplifications of this worked run: $\bar s_{\rm wet}=0.6$ and $\kappa=1$ in the root advance; no spacing rule.)}\sol \givenline{$\bar I=150$ mm yr$^{-1}$ find ; $P=600$ mm ; $E_{\rm yr}=170$ ; $T_{\rm yr}=230$ ; $k_t=1$ ; $\bar s_{\rm wet}=0.6$ ; $\kappa=1$}\ans Dose: $150/10=15$ mm, which is at the lower clip, so $d=15$ mm. (The run below used the limit rule $\mathrm{clip}(15;15,35)$.) The run gives: rain $P=606$ mm, irrigation $90$ mm in 6 events on days 147, 148, 164, 165, 181, 198, soil evaporation $158$ mm, transpiration $230$ mm, drainage $200$ mm, and final rooting depth $82$ cm. Storage change: the profile started at $0.72W_{\max}=3.06$ mm per layer, 122.4 mm in all. \finans

\subhd{Balance check.} $P+I-E-T-D=\Delta S$: $606+90-158-230-200=108$ mm; the throughfall is smaller than $P$ by the interception $P-\sum T_f$, so the storage change equals this minus the interception of 62 mm. In words: the engine conserves mass because every flux is taken from, or given to, the layers.\par
The first irrigation occurs on day 147, when the sensor layer first reaches $1-s_8\ge0.468$; before that the crop lived on rain and storage.


\exbox{The same season in the sand}{Repeat the season run for the sand S03 (AWC 0.6 mm cm$^{-1}$, trigger depletion $f_T=0.842$ at 300 cm) with the same weather seed.}\sol \givenline{$f_T=0.842$}\ans $W_{\max}=0.6\times2.5=1.5$ mm per layer. Irrigation $90$ mm in 6 events (days 140, 155, 165, 174, 188, 199), drainage $251$ mm, transpiration $226$ mm, evaporation $142$ mm. Compared with the loam (drainage $200$ mm) the sand drains $51$ mm more: the layers hold only 1.5 mm each, so most of every shower and irrigation passes below 1 m. With the same dose of 15 mm the sand wastes a larger share, which is the mechanism behind the fate of saved water in Chapter~\ref{ch:fate}. \finans

\subhd{Reading.} The daily model is a tipping-bucket engine; the application scales drainage and the wetting front to Engine~1 (Equation below) because they are not Richards solutions.


\section{Scaling to the annual model}
The daily results are multiplied by constants so that their 20-year means equal Engine 1:
\begin{equation}\boxed{S_P=\frac P{P_{\rm sim}},\ S_{\rm thr}=\frac{P-P_{\rm other}}{T_{f,\rm sim}},\ S_E=\frac E{E_{\rm sim}},\ S_T=\frac T{T_{\rm sim}},\ S_D=\frac D{D_{\rm sim}},}\end{equation}
\par\noindent where $S_P,S_E,S_T,S_D$ are the scale factors of rain, evaporation, transpiration and drainage (dimensionless), $P$, $E$, $T$, $D$ the target values and the subscript ``sim'' the simulated values (mm), $P_{\rm other}$ the rain not intercepted, and $T_{f,{\rm sim}}$ the simulated throughfall.

and the displayed fluxes are $R=P_dS_P$, $T_f=\min\{R,T_{f,d}S_{\rm thr}\}$, $E=E_dS_E$, $T=T_dS_T$, $D=D_dS_D$, with $\Delta S_d=T_f+I-E-T-D$.

\exbox{Scale factors}{Engine 1 gives $P=600$, $P_{\rm other}=60$, $E=170$, $T=230$, $D=70$ mm. The daily run of the loam above gave $P_{\rm sim}=606.1$, $T_{f,\rm sim}=544.4$, $E_{\rm sim}=158.5$, $T_{\rm sim}=229.6$, $D_{\rm sim}=200.2$ mm. Find the scale factors and the displayed drainage of a day on which the model drained $12$ mm.}\sol \givenline{$P=600$ ; $P_{\rm other}=60$ ; $E=170$ ; $T=230$ ; $D=70$ mm ; $P_{\rm sim}=606.1$ ; $T_{f,\rm sim}=544.4$ ; $E_{\rm sim}=158.5$ ; $T_{\rm sim}=229.6$}\ans $S_P=600/606.1=0.9899$; $S_{\rm thr}=540/544.4=0.9919$; $S_E=170/158.5=1.0726$; $S_T=230/229.6=1.0018$; $S_D=70/200.2=0.3496$. A day with a drainage of 12 mm is shown as $12\times0.3496=4.20$ mm. The scale factors are close to 1 when the daily engine is already consistent with the annual one; a factor far from 1 (here $S_D=0.35$) flags a bucket model that drains too much compared with the Richards-based Daisy result, which is the reason for the scaling. \finans


\exbox{Evaporation demand under warming}{Under warming the bare-soil demand is $E_{0,\rm bare}=4.2(1+0.07\Delta T)\cdot\frac{0.35+0.65\cdot\frac12[1+\sin(2\pi(\mathrm{doy}-100)/365)]}{0.675}\,\mathrm{clip}(1-0.12(\rho-1);0.85,1.12)\sqrt{\gamma_y}$ mm d$^{-1}$. Find the demand on day 190 with $\Delta T=2$ K, rain scale $\rho=1$, $\gamma_y=1$.}\sol \givenline{$\Delta T=2$ K ; $\rho=1$ ; $\gamma_y=1$}\ans $4.2\times1.14=4.788$; the seasonal weight is $0.35+0.325[1+\sin(2\pi\,90/365)]=0.35+0.325(1+0.9998)=0.9999$, divided by $0.675$ gives $1.4814$; the cloud factor is $\mathrm{clip}(1,0.85,1.12)=1$ and $\sqrt{\gamma_y}=1$. So $E_{0,\rm bare}=7.09$ mm d$^{-1}$, which is 14 per cent above the unwarmed value 6.22 mm d$^{-1}$ (a 7 per cent per kelvin increase, close to the Clausius--Clapeyron rate).
\banner{Field balance, soil--plant--air and profile evolution}\label{ch:bal}
\section{Field balance: the closed budget}\label{sec:closure}
The home layer draws the water budget of a field for a day, a year or the 20-year mean: $P+I=T+E+E_{\rm other}+D+\Delta S$ (Equation 6.12 of the application). In the average-year view the five numbers over the arrows and the strip below must satisfy this identity. \finans

\exbox{Checking the budget and changing the rain}{The average-year view shows $P=598$, $I=142$, $T=236$, $E=171$, interception and mulch evaporation $E_{\rm other}=88$, $D=207$ mm. (These are illustrative numbers for the check.) (a) Compute the storage change. (b) Change the rainfall scale to $\times1.3$ and recompute the rain; if the extra rain is shared as $25$ per cent evaporated, $15$ per cent intercepted and the rest drained, find the new drainage and $\Delta S$.}\sol \givenline{$P=598$ ; $I=142$ ; $T=236$ ; $E=171$ ; $E_{\rm other}=88$ ; $D=207$ mm}\ans (a) Inputs $598+142=740$ mm; outputs $236+171+88+207=702$ mm; $\Delta S=740-702=+38$ mm. (b) The new rain is $1.3\times598=777.4$ mm, an extra $179.4$ mm. Evaporation takes $0.25\times179.4=44.9$ mm, interception $0.15\times179.4=26.9$ mm, drainage the remaining $60$ per cent, $107.6$ mm. New values: $E=215.9$, $E_{\rm other}=114.9$, $D=314.6$, and $\Delta S=777.4+142-(236+215.9+114.9+314.6)=+38$ mm (unchanged by construction). Drainage rises by 52 per cent for 30 per cent more rain, which is why the drainage and nitrate-leaching arrows thicken together. \finans

\subhd{Efficiency indices.} The irrigation efficiency for a yield of 6.6 Mg ha$^{-1}$ is $100\times6.6/142=4.65$ kg m$^{-3}$; the transpiration share of evapotranspiration is $236/(236+171+88)=47.7$ per cent.


\section{Soil, plant and air: potentials along the path}
The suction of each layer is $h_i=\max\{30,h(f_i)\}$ with $f_i=1-s_i$ and $\psi_{\rm soil,i}=-h_i\times9.807\times10^{-5}$ MPa cm$^{-1}$. The root-zone suction is the density-weighted mean $h_{RZ}=\sum\rho_ih_i/\sum\rho_i$. The plant is a chain of series resistances,
\begin{align}
R_{\rm soil}&=0.03\max\{K_r,0.002\}^{-0.55},\quad R_{\rm root}=0.10,\ R_{\rm xyl}=0.05,\ R_{\rm leaf}=0.08\ \ (\text{MPa per mm d}^{-1}),\\
\psi_{\rm root}&=\psi_{RZ}-aER_{\rm soil},\quad\psi_{\rm stem}=\psi_{\rm root}-aER_{\rm root},\quad\psi_{\rm leaf}=\mathrm{clip}\bigl(\psi_{\rm stem}-aE(R_{\rm xyl}+R_{\rm leaf});-3.2,0\bigr),
\end{align}
with $a=1$ when the cover exceeds 0.05 and $a=0$ for a fallow field. The air above the canopy has $\psi_{\rm air}=(RT/V_w)\ln\mathrm{RH}\approx135.5\ln\mathrm{RH}$ MPa.

\exbox{The soil--plant--air path}{A transpiring crop ($a=1$) takes $E=4$ mm d$^{-1}$. Compute the potentials along the path for the loam S06 at root-zone suctions of 50, 300 and 1000 cm, and for the sand S03 at 300 cm.}
\textbf{Solution (loam, 300 cm, in full).} 1. Relative conductivity $K_r(300)=3.8\times10^{-5}$, below the floor $0.002$, so $R_{\rm soil}=0.03(0.002)^{-0.55}=0.03\times30.5=0.915$ MPa per mm d$^{-1}$. 2. Soil potential $\psi_{RZ}=-300\times9.807\times10^{-5}=-0.0294$ MPa. 3. $\psi_{\rm root}=-0.0294-4(0.915)=-3.689$ MPa. 4. $\psi_{\rm stem}=-3.689-4(0.10)=-4.089$ MPa. 5. $\psi_{\rm leaf}=-4.089-4(0.05+0.08)=-4.609$, clipped to $-3.2$ MPa. Table~\ref{tab:spa} lists the other cases.\begin{table}[!htbp]\centering\small\caption{Potentials (MPa) along the soil--plant path for $E=4$ mm d$^{-1}$.}\label{tab:spa}
\begin{tabular}{lrrrrrrr}\toprule
\hd{Soil} & \hd{$h$ (cm)} & \hd{$K_r$} & \hd{$R_{\rm soil}$} & \hd{$\psi_{RZ}$} & \hd{$\psi_{\rm root}$} & \hd{$\psi_{\rm stem}$} & \hd{$\psi_{\rm leaf}$}\\\midrule
S06 & 50 & 1.03e-02 & 0.371 & -0.0049 & -1.489 & -1.889 & -2.409\\
S06 & 300 & 3.80e-05 & 0.915 & -0.0294 & -3.691 & -4.091 & -3.200\\
S06 & 1000 & 6.55e-07 & 0.915 & -0.0981 & -3.759 & -4.159 & -3.200\\
S03 & 300 & 2.73e-11 & 0.915 & -0.0294 & -3.691 & -4.091 & -3.200\\\bottomrule\end{tabular}\end{table}
\textbf{Caution (reproduced by the numbers).} In every case except the wet loam the relative conductivity is below the floor and $R_{\rm soil}$ sits at its cap of $0.92$ MPa per mm d$^{-1}$; the leaf potential then reaches the clip at $-3.2$ MPa, below the intended range $-1$ to $-1.5$ MPa. The leaf, stem and root potentials of this layer therefore illustrate the series-resistance concept and are not a measure of plant water status.


\exbox{The water potential of the air}{Compute the air potential at $\mathrm{RH}=66$ per cent and under a warming of $\Delta T=3$ K, where $\mathrm{RH}=\mathrm{clip}(0.66-0.03\Delta T;0.35,0.90)$. Also verify the constant 135.5.}\sol \givenline{$\mathrm{RH}=66$ per cent ; $\Delta T=3$ K ; $\mathrm{RH}=\mathrm{clip}(0.66-0.03\Delta T;0.35,0.90)$}\ans $RT/V_w=8.314\times293/(18\times10^{-6})=1.353\times10^{8}$ Pa $=135.3$ MPa (the application uses 135.5). At $66$ per cent: $\psi_{\rm air}=135.5\ln0.66=-56.3$ MPa. With $\Delta T=3$ K the humidity is $0.66-0.09=0.57$ and $\psi_{\rm air}=135.5\ln0.57=-76.2$ MPa: the air gets $20$ MPa drier. The potential difference from a leaf at $-1.5$ MPa to the air is therefore about 55 to 75 MPa, far larger than any difference inside the plant; the resistances of the chain decide how much of it is used, not the size of the air potential itself. \finans


\section{Profile evolution}
The share of the 0 to 60 cm root zone beyond the trigger is the thickness-weighted fraction
\begin{equation}\boxed{\mathrm{share}=\frac{\sum_{i:z_i\le60,\ h_i>h_T}\Delta z_i}{\sum_{i:z_i\le60}\Delta z_i},}\end{equation}
\par\noindent where The share is the fraction of the top 60 cm of the profile whose suction $h_i$ exceeds the trigger $h_T$, $\Delta z_i$ is the thickness of layer $i$ (cm) and $z_i$ its depth (cm).

and the seasonal mean suction is the geometric mean $m_y=10^{\frac1{183}\sum_d\log_{10}h_d}$.

\exbox{Share of the root zone beyond the trigger}{Twelve layers of 5 cm (down to 60 cm) have suctions $h_i=60,80,120,200,350,500,700,900,1100,1300,1600,2000$ cm. Find the share beyond a trigger of 300 cm and beyond 1000 cm; and the geometric mean suction of the profile.}\sol \givenline{$h_i=60,80,120,200,350,500,700,900,1100,1300,1600,2000$ cm}\ans Layers above 300 cm: 350, 500, 700, 900, 1100, 1300, 1600, 2000, which is 8 of 12, so the share is $8\times5/60=0.667$. Beyond 1000 cm: 1100, 1300, 1600, 2000, which is 4 of 12, so $0.333$. The geometric mean is $10^{\overline{\log_{10}h}}$: the mean of the logarithms is $2.6470$, so $m=444$ cm, which is lower than the arithmetic mean $742$ cm because the geometric mean damps the very dry layers. The heat strip of the application uses the geometric mean for this reason. \finans


\exbox{Dry seasons in the record}{In the 20-year record the median growing-season mean suction is $m_{\rm med}=250$ cm. A season has $\overline{\log_{10}h}=2.8$. Find its $m_y$, the year factor $\gamma_y$ for rainfall generation, and the mean wet-day depth for $P=600$ mm.}\sol \givenline{$m_{\rm med}=250$ cm ; $\overline{\log_{10}h}=2.8$ ; $P=600$ mm}\ans $m_y=10^{2.8}=631$ cm. $\gamma_y=(250/631)^{0.3}=0.396^{0.3}=e^{0.3\ln0.396}=e^{-0.2788}=0.757$ (inside $[0.65,1.45]$). $\mu_0=5.871$ mm, so the mean wet-day depth is $5.871\times0.757=4.44$ mm and the expected annual rain is $0.28\times365\times4.44=454$ mm, 24 per cent below the 600 mm normal. A dry season in the record gives a dry generated year.
\banner{The Dashboard: how every number is produced}\label{ch:dash}
The Dashboard recomputes from the engine every time the state, the date or the weather changes, so it adds no model of its own: every number is a read-out or a sum of the daily series $P_k,I_k,E_k,T_k,D_k$ of Engine~2 (Chapter~\ref{ch:e2}) and of the annual outputs of Engine~1. This chapter does the sums. \finans

\section{Block 1: the working date}
On day index $k$ the tab shows $\mathrm{Rain}=P_k$, $\mathrm{Irr}=I_k$, $\mathrm{Evap}=E_k$, $\mathrm{Transp}=T_k$, $\mathrm{Drain}=D_k$ (mm d$^{-1}$), the sensor depletion $f$ and the trigger $f_T$ (irrigation fires when $f\ge f_T$).

\exbox{One day of the dashboard}{The loam at the trigger $-300$ cm (depletion $f_T=0.468$), the first synthetic year (seed 1), working date day 150. Read the first block.}\sol \givenline{$f_T=0.468$}\ans From the daily series: rain $P=1.0$ mm; rain reaching the soil after interception $0.3$ mm; irrigation $I=0.0$ mm; soil evaporation $E=0.54$ mm; transpiration $T=2.20$ mm; drainage $D=0.0$ mm; canopy cover $c=0.48$; the sensor layer holds 0.924 of its capacity so its relative depletion is $f=1-0.924=0.076$, below the trigger $f_T=0.468$. Irrigation fires on a day when the sensor depletion has reached $f_T$ (Eq.~6.1). \finans


\section{Block 2: the next 14 days}
The sums over the following $N=14$ days are $\Sigma X=\sum_{j=1}^{14}X_{k+j}$ for $X\in\{P,I,E,T,D\}$. The bar says whether the days come from the weather record (the analogue year) or from the 16-day forecast when online weather is in use.

\exbox{The next fortnight}{Sum days 151 to 164.}\sol  $\Sigma P=13.6$ mm, $\Sigma I=15.0$ mm, $\Sigma E=7.4$ mm, $\Sigma T=39.8$ mm, $\Sigma D=0.0$ mm. Evapotranspiration of the fortnight is $\Sigma E+\Sigma T=47.2$ mm, a mean of $3.37$ mm d$^{-1}$. Without rain or irrigation the fortnight would draw 18.6 mm from storage (the model gives the actual change below). \finans


\section{Block 3: the whole year and its closure}
For the calendar year the tab gives rain, irrigation, other evaporation (interception, litter, ponded water, snow), soil evaporation, transpiration, drainage and the storage change; the balance closes by the identity
\begin{equation}\boxed{\Delta S=P+I-E_{\rm other}-E-T-D.}\end{equation}
\par\noindent where $\Delta S$ is the change of stored water, $P$ the rain, $I$ the irrigation, $E_{\rm other}$ the interception and surface evaporation, $E$ the soil evaporation, $T$ the transpiration and $D$ the drainage (all in mm).


\begin{table}[!htbp]\centering\small\caption{Block 3 of the Results tab.}\label{tab:year}
\begin{tabular}{lr}\toprule
\hd{Quantity} & \hd{mm (year 1, loam, $-300$ cm)}\\\midrule
Rain $P$ & 606.1\\
Irrigation $I$ & 90.0\\
Other evaporation $E_{\rm other}$ & 61.7\\
Soil evaporation $E$ & 158.5\\
Transpiration $T$ & 229.6\\
Drainage $D$ & 200.2\\
Storage change $\Delta S$ & 46.1\\
Deepest roots (cm) & 82\\\bottomrule\end{tabular}\end{table}

\exbox{Closing the year}{Verify the closure of Table~\ref{tab:year}.}\sol  $\Delta S=P+I-E_{\rm other}-E-T-D=606.1+90.0-61.7-158.5-229.6-200.2=46.1$ mm. Independently, the root-zone store of the 100 cm profile rose from 122.4 mm on 1 January to 168.5 mm on 31 December, a change of 46.1 mm. The two agree because $E_{\rm other}$ is, by construction, the rain that did not reach the soil: $E_{\rm other}=\sum(P_k-P_k^{\rm soil})=61.7$ mm. \finans


\section{Block 4 and the Charts tab: twenty seasons}
Block 4 shows the 20-year averages of the annual model, and each chart has a line stating the mean, the 10--90\% range, the lowest and highest season and the variation between years. The variation is the coefficient of variation $\mathrm{CV}=100\,s/\bar x$ with $s$ the sample standard deviation.

\begin{table}[!htbp]\centering\small\caption{Reliability over twenty synthetic seasons (seeds 1--20), loam, trigger $-300$ cm, mm per year.}\label{tab:rel}
\begin{tabular}{lrrrrrr}\toprule
\hd{Quantity} & \hd{mean} & \hd{10\%} & \hd{90\%} & \hd{lowest} & \hd{highest} & \hd{CV (\%)}\\\midrule
Rain & 602 & 522 & 735 & 438 & 765 & 14\\
Irrigation & 99 & 58 & 137 & 45 & 150 & 32\\
Drainage & 207 & 134 & 283 & 111 & 359 & 31\\
Soil evaporation & 158 & 157 & 160 & 148 & 160 & 2\\
Transpiration & 230 & 229 & 230 & 228 & 230 & 0\\\bottomrule\end{tabular}\end{table}

\exbox{The statistics line}{From Table~\ref{tab:rel} write the line under the irrigation chart.}\sol  ``Irrigation: mean 99 mm; 10 to 90\% range 58 to 137 mm; lowest 45 mm, highest 150 mm; variation between years 32\%.'' The coefficient of variation is $100\times s/\bar x$ with $s=31.3$ mm and $\bar x=99$ mm. No season was rainfed in this sample. \finans


\section{The pinned baseline}
Press Pin and the tiles of Block 4 are stored; afterwards each tile shows $\Delta_i=x_i-x_i^{\rm pin}$ with $\blacktriangle$ for an increase, $\blacktriangledown$ for a decrease, green when the change is in the better direction, red when worse and ``= no change'' when $|\Delta_i|<10^{-9}$. Lower is better for soil evaporation, drainage, irrigation and nitrogen leaching; higher is better for yield and irrigation efficiency; transpiration, evapotranspiration and rainfall are neutral.

\begin{table}[!htbp]\centering\small\caption{Pin at the trigger $-300$ cm, then move to $-1500$ cm (twenty-season means, mm per year).}\label{tab:pin}
\begin{tabular}{lrrrll}\toprule
\hd{Tile} & \hd{pinned (-300 cm)} & \hd{now (-1500 cm)} & \hd{$\Delta$} & \hd{arrow} & \hd{colour}\\\midrule
Soil evaporation & 157.8 & 156.1 & -1.7 & $\blacktriangledown$ & better (green)\\
Transpiration & 229.6 & 229.6 & 0.0 & = no change & neutral\\
Drainage & 206.9 & 156.2 & -50.7 & $\blacktriangledown$ & better (green)\\
Irrigation & 99.0 & 45.8 & -53.2 & $\blacktriangledown$ & better (green)\\
Rain & 602.0 & 602.0 & 0.0 & = no change & neutral\\\bottomrule\end{tabular}\end{table}

\exbox{Reading the pinned differences}{Read the tiles of Table~\ref{tab:pin}.}\sol  Irrigation falls from 99.0 to 45.8 mm (-53.2 mm): $\blacktriangledown$ and green, since less is better. Transpiration changes by 0.0 mm and is neutral (it carries no colour, because less transpiration may mean either less stress or a smaller crop). Drainage changes by -50.7 mm. A drier trigger therefore saves irrigation but the dashboard alone cannot say whether the saving is worth it: the yield tile (not computed here) must be read with it, as in Chapter~\ref{ch:ctrl}. \finans


\section{The pump hydrograph of a design}
For drip and sprinkler designs the pump hydrograph is the step series $Q_{\rm pump}(t)=Q_z$ for $t\in[0,n_zt_{\rm run})$ and 0 afterwards: the zones or sets run in turn.

\exbox{A pump hydrograph}{For the sprinkler farm of Chapter~\ref{ch:sprinkler} ($n=3$ laterals at once, set time $T=371$ min) the farm has $N_m=18$ lateral positions, so $n_z=\lceil18/3\rceil=6$ sets. Find the step series.}\sol \givenline{$n=3$ laterals ; $T=371$ min ; $N_m=18$ lateral positions ; $n_z=\lceil18/3\rceil=6$ sets}\ans The flow of one set is $Q_z=3\times12\times21.2$ L min$^{-1}=764$ L min$^{-1}=12.74$ L s$^{-1}$ (3 laterals of 12 sprinklers). $t_{\rm run}=371/60=6.18$ h. The pump runs at $Q_z$ for $n_zt_{\rm run}=37.1$ h, in steps of $6.18$ h, and the volume is $Q_z\times n_zt_{\rm run}\times3.6=1{,}701$ m$^3$ (check: gross depth $\times$ area of the sets). \finans


\section{Summary}
Every Dashboard number is a read-out, a sum, a mean or a difference of the daily series: Block 1 reads day $k$; Block 2 sums 14 days; Block 3 closes $\Delta S=P+I-E_{\rm other}-E-T-D$; Block 4 and the charts use twenty seasons; the pin shows $\Delta=x-x^{\rm pin}$ with a colour by direction; the pump hydrograph is a step series. Because all of it comes from the same objects, the dashboard, the report and the CSV exports agree digit for digit.
\IfFileExists{extra/app08.tex}{\subsection{The output of the application, panel by panel}
The engines of Chapter 9 are controlled from the dashboard of the application, and the dashboard shows their results in four tabs: the \emph{Results} table, the \emph{Design} picture, the \emph{Charts} and the \emph{Messages} of the run.

The picture shows the soil, the crop, the mulch and the fluxes of the season; the numbers at the right give the irrigation, the drainage, the evaporation and the yield.

\appfig{dash_results}{1.0}{The Dashboard, Results tab: the picture of the field with the water fluxes and the numbers of the run.}

The charts show the daily balance of the bucket engine of this chapter: rain, irrigation, evapotranspiration and drainage with the depletion against the trigger.

\appfig{dash_charts}{1.0}{The Dashboard, Charts tab: the daily water balance and the depletion of the root zone.}

The Design tab shows the layout of the selected design with the numbers of the engine.

\appfig{dash_design}{1.0}{The Dashboard, Design tab: the design of the selected irrigation layer.}

Every step of the run is written in the message log, so that the user can see which engine ran and what it found.

\appfig{dash_messages}{1.0}{The Dashboard, Messages tab: the log of the run, with the warnings of the engine.}

}{}
\IfFileExists{extra/long08.tex}{\subsection{Ten days of the bucket engine}
\exbox{Daily bucket engine step by step}{Run the daily bucket engine for ten days with $TAW=140$ mm, $p=0.5$, $K_c=1.1$, the given $ET_0$ and rain, an initial depletion of 40 mm and a rule that irrigates 40 mm whenever the depletion passes $RAW$.}\sol
\ans Initial depletion $D_0=40$ mm, $TAW=140$ mm, $RAW=pTAW=70$ mm, $K_c=1.1$; irrigation of 40 mm is added when $D_r>RAW$
\ans Day 1: $D=40.0$ mm, $K_s=1.00$, $ET_a=K_sK_cET_0=5.72$, $P=0$ mm; irrigation 0 mm; $D_{new}=45.7$ mm
\ans Day 2: $D=45.7$ mm, $K_s=1.00$, $ET_a=K_sK_cET_0=6.16$, $P=0$ mm; irrigation 0 mm; $D_{new}=51.9$ mm
\ans Day 3: $D=51.9$ mm, $K_s=1.00$, $ET_a=K_sK_cET_0=6.60$, $P=0$ mm; irrigation 0 mm; $D_{new}=58.5$ mm
\ans Day 4: $D=58.5$ mm, $K_s=1.00$, $ET_a=K_sK_cET_0=6.38$, $P=0$ mm; irrigation 0 mm; $D_{new}=64.9$ mm
\ans Day 5: $D=64.9$ mm, $K_s=1.00$, $ET_a=K_sK_cET_0=6.71$, $P=0$ mm; irrigation 40 mm; $D_{new}=31.6$ mm
\ans Day 6: $D=31.6$ mm, $K_s=1.00$, $ET_a=K_sK_cET_0=7.04$, $P=0$ mm; irrigation 0 mm; $D_{new}=38.6$ mm
\ans Day 7: $D=38.6$ mm, $K_s=1.00$, $ET_a=K_sK_cET_0=6.49$, $P=12$ mm; irrigation 0 mm; $D_{new}=33.1$ mm
\ans Day 8: $D=33.1$ mm, $K_s=1.00$, $ET_a=K_sK_cET_0=5.28$, $P=3$ mm; irrigation 0 mm; $D_{new}=35.4$ mm
\ans Day 9: $D=35.4$ mm, $K_s=1.00$, $ET_a=K_sK_cET_0=5.50$, $P=0$ mm; irrigation 0 mm; $D_{new}=40.9$ mm
\ans Day 10: $D=40.9$ mm, $K_s=1.00$, $ET_a=K_sK_cET_0=6.05$, $P=0$ mm; irrigation 0 mm; $D_{new}=46.9$ mm


}{}
\IfFileExists{extra/proj08.tex}{\banner{Design project 8: choosing the trigger of a maize season by simulation}\label{sec:proj8}
\textbf{Brief.} A season of 125 days of maize on a soil with TAW = 150 mm and $p=0.5$ is simulated with the daily bucket on a weather record of seeded rain events (a mean of 212 mm of rain, $\Sigma ET_0=797$ mm). The irrigation rule applies the whole depletion when it reaches a fraction $f_t$ of TAW. The designer compares triggers from $f_t=0.3$ to 0.8 and the rainfed case on irrigation, drainage, yield and the extra yield per extra millimetre of irrigation. The maximum yield is $Y_m=11.5$ Mg\,ha$^{-1}$ and $K_y=1.28$.

\begin{table}[!htbp]\centering\small\caption{Seasonal results of the bucket for seven triggers.}
\begin{tabular}{lrrrrrr}\toprule
\hd{Trigger $f_t$} & \hd{Irrigations} & \hd{Irrigation (mm)} & \hd{Drainage (mm)} & \hd{$ET_a$ (mm)} & \hd{$Y/Y_m$} & \hd{Yield gain per 100 mm}\\\midrule
0.30 & 12 & 577 & 45 & 747 & 1.000 & 2.0\\
0.40 & 9 & 573 & 41 & 747 & 1.000 & 2.0\\
0.50 & 7 & 556 & 30 & 747 & 1.000 & 2.1\\
0.60 & 6 & 557 & 57 & 737 & 0.983 & 2.0\\
0.70 & 4 & 430 & 30 & 699 & 0.919 & 2.5\\
0.80 & 3 & 367 & 30 & 641 & 0.818 & 2.6\\
rainfed & 0 & 0 & 30 & 292 & 0.221 & 0.0\\\bottomrule\end{tabular}\end{table}

\textbf{Reading the table.} A trigger below the stress threshold ($f_t\le0.5=p$) keeps $K_s=1$, so $ET_a=ET_c$ and the relative yield is 1; lower triggers add irrigations but not yield and send water to drainage. Triggers above $p$ let stress start and the yield fall by $K_y$ times the relative deficit. The rainfed season gives $Y/Y_m=0.221$. The choice of the designer is $f_t=0.5$: the largest depletion that keeps the yield at its maximum, which needs the least irrigation (556 mm) among the stress-free triggers.

\textbf{Check.} A trigger of 0.6 saves -1 mm of irrigation against 0.5 at a yield loss of 1.7 per cent; whether it is worth it depends on the price of water against the price of grain, as in Chapter~5.


\fig{g8_bucket}{0.95}{Daily bucket for a 125-day maize season. (a) Depletion for three triggers, with the irrigation applied in brackets and the limit RAW. (b) The rainfall of the season.}
}{}
\IfFileExists{extra/ten08.tex}{\subsection{Worked solutions on the ten soils}
Runoff by the curve number and the fate of the rain are computed for a storm of 60 mm on each soil.

\exbox{Ten-soil example 8a: Coarse loam (S01)}{Coarse loam (hydrologic group B, $CN=75$): a storm of 60 mm falls on the soil. Find the runoff, the effective rainfall and the share that fills the bucket of $TAW=90$ mm if the profile was at 50 per cent depletion.}\sol \givenline{$CN=75$ ; $TAW=90$ mm}\ans $S=25\,400/75-254=84.7$ mm; $I_a=0.2S=16.9$ mm \ans \[Q=(60-16.9)^2/(60-16.9+84.7)=14.5\ \text{mm}\] \ans Effective rain $60-14.5=45.5$ mm \ans The bucket holds $RAW=45$ mm of room, so it overfills: drainage 0.5 mm \finans

\exbox{Ten-soil example 8a: Coarse sand (S02)}{Coarse sand (hydrologic group A, $CN=65$): a storm of 60 mm falls on the soil. Find the runoff, the effective rainfall and the share that fills the bucket of $TAW=30$ mm if the profile was at 50 per cent depletion.}\sol \givenline{$CN=65$ ; $TAW=30$ mm}\ans $S=25\,400/65-254=136.8$ mm; $I_a=0.2S=27.4$ mm \ans \[Q=(60-27.4)^2/(60-27.4+136.8)=6.3\ \text{mm}\] \ans Effective rain $60-6.3=53.7$ mm \ans The bucket holds $RAW=15$ mm of room, so it overfills: drainage 38.7 mm \finans

\exbox{Ten-soil example 8a: Sand (S03)}{Sand (hydrologic group A, $CN=65$): a storm of 60 mm falls on the soil. Find the runoff, the effective rainfall and the share that fills the bucket of $TAW=36$ mm if the profile was at 50 per cent depletion.}\sol \givenline{$CN=65$ ; $TAW=36$ mm}\ans $S=25\,400/65-254=136.8$ mm; $I_a=0.2S=27.4$ mm \ans \[Q=(60-27.4)^2/(60-27.4+136.8)=6.3\ \text{mm}\] \ans Effective rain $60-6.3=53.7$ mm \ans The bucket holds $RAW=18$ mm of room, so it overfills: drainage 35.7 mm \finans

\exbox{Ten-soil example 8a: Loamy sand (S04)}{Loamy sand (hydrologic group A, $CN=65$): a storm of 60 mm falls on the soil. Find the runoff, the effective rainfall and the share that fills the bucket of $TAW=54$ mm if the profile was at 50 per cent depletion.}\sol \givenline{$CN=65$ ; $TAW=54$ mm}\ans $S=25\,400/65-254=136.8$ mm; $I_a=0.2S=27.4$ mm \ans \[Q=(60-27.4)^2/(60-27.4+136.8)=6.3\ \text{mm}\] \ans Effective rain $60-6.3=53.7$ mm \ans The bucket holds $RAW=27$ mm of room, so it overfills: drainage 26.7 mm \finans

\exbox{Ten-soil example 8a: Sandy loam (S05)}{Sandy loam (hydrologic group B, $CN=75$): a storm of 60 mm falls on the soil. Find the runoff, the effective rainfall and the share that fills the bucket of $TAW=72$ mm if the profile was at 50 per cent depletion.}\sol \givenline{$CN=75$ ; $TAW=72$ mm}\ans $S=25\,400/75-254=84.7$ mm; $I_a=0.2S=16.9$ mm \ans \[Q=(60-16.9)^2/(60-16.9+84.7)=14.5\ \text{mm}\] \ans Effective rain $60-14.5=45.5$ mm \ans The bucket holds $RAW=36$ mm of room, so it overfills: drainage 9.5 mm \finans

\exbox{Ten-soil example 8a: Loam (S06)}{Loam (hydrologic group B, $CN=75$): a storm of 60 mm falls on the soil. Find the runoff, the effective rainfall and the share that fills the bucket of $TAW=102$ mm if the profile was at 50 per cent depletion.}\sol \givenline{$CN=75$ ; $TAW=102$ mm}\ans $S=25\,400/75-254=84.7$ mm; $I_a=0.2S=16.9$ mm \ans \[Q=(60-16.9)^2/(60-16.9+84.7)=14.5\ \text{mm}\] \ans Effective rain $60-14.5=45.5$ mm \ans The bucket holds $RAW=51$ mm of room, so it absorbs all: drainage 0.0 mm \finans

\exbox{Ten-soil example 8a: Silt loam (S07)}{Silt loam (hydrologic group C, $CN=82$): a storm of 60 mm falls on the soil. Find the runoff, the effective rainfall and the share that fills the bucket of $TAW=120$ mm if the profile was at 50 per cent depletion.}\sol \givenline{$CN=82$ ; $TAW=120$ mm}\ans $S=25\,400/82-254=55.8$ mm; $I_a=0.2S=11.2$ mm \ans \[Q=(60-11.2)^2/(60-11.2+55.8)=22.8\ \text{mm}\] \ans Effective rain $60-22.8=37.2$ mm \ans The bucket holds $RAW=60$ mm of room, so it absorbs all: drainage 0.0 mm \finans

\exbox{Ten-soil example 8a: Sandy clay loam (S08)}{Sandy clay loam (hydrologic group C, $CN=82$): a storm of 60 mm falls on the soil. Find the runoff, the effective rainfall and the share that fills the bucket of $TAW=84$ mm if the profile was at 50 per cent depletion.}\sol \givenline{$CN=82$ ; $TAW=84$ mm}\ans $S=25\,400/82-254=55.8$ mm; $I_a=0.2S=11.2$ mm \ans \[Q=(60-11.2)^2/(60-11.2+55.8)=22.8\ \text{mm}\] \ans Effective rain $60-22.8=37.2$ mm \ans The bucket holds $RAW=42$ mm of room, so it absorbs all: drainage 0.0 mm \finans

\exbox{Ten-soil example 8a: Clay loam (S09)}{Clay loam (hydrologic group D, $CN=87$): a storm of 60 mm falls on the soil. Find the runoff, the effective rainfall and the share that fills the bucket of $TAW=96$ mm if the profile was at 50 per cent depletion.}\sol \givenline{$CN=87$ ; $TAW=96$ mm}\ans $S=25\,400/87-254=38.0$ mm; $I_a=0.2S=7.6$ mm \ans \[Q=(60-7.6)^2/(60-7.6+38.0)=30.4\ \text{mm}\] \ans Effective rain $60-30.4=29.6$ mm \ans The bucket holds $RAW=48$ mm of room, so it absorbs all: drainage 0.0 mm \finans

\exbox{Ten-soil example 8a: Silty clay loam (S10)}{Silty clay loam (hydrologic group D, $CN=87$): a storm of 60 mm falls on the soil. Find the runoff, the effective rainfall and the share that fills the bucket of $TAW=102$ mm if the profile was at 50 per cent depletion.}\sol \givenline{$CN=87$ ; $TAW=102$ mm}\ans $S=25\,400/87-254=38.0$ mm; $I_a=0.2S=7.6$ mm \ans \[Q=(60-7.6)^2/(60-7.6+38.0)=30.4\ \text{mm}\] \ans Effective rain $60-30.4=29.6$ mm \ans The bucket holds $RAW=51$ mm of room, so it absorbs all: drainage 0.0 mm \finans

}{}
\IfFileExists{extra/tenb08.tex}{\subsection{Worked solutions on ten crops}
The ten crops of the library differ in their demand, their roots and their sensitivity. The same type of question is answered for each of them at $ET_0=5.5$ mm\,d$^{-1}$; the data are typical values of the crops (Doorenbos--Kassam and FAO 56).

\exbox{Crop example 8a: Maize}{Maize: a full bucket ($D=0$) drains by evapotranspiration at $K_cET_0=6.60$ mm\,d$^{-1}$ with no rain. After how many days does $K_s$ fall below 1, and what are $D$ and $K_s$ five days later?}\sol $TAW=170$, $RAW=94$ mm \ans $D$ reaches RAW after $14.2$ d \ans Then each day $D\leftarrow D+K_sK_cET_0$ with $K_s=(TAW-D)/(TAW-RAW)$; after five days $D=121.3$ mm \ans $K_s=0.64$ \ans

\exbox{Crop example 8a: Wheat}{Wheat: a full bucket ($D=0$) drains by evapotranspiration at $K_cET_0=6.32$ mm\,d$^{-1}$ with no rain. After how many days does $K_s$ fall below 1, and what are $D$ and $K_s$ five days later?}\sol $TAW=204$, $RAW=112$ mm \ans $D$ reaches RAW after $17.7$ d \ans Then each day $D\leftarrow D+K_sK_cET_0$ with $K_s=(TAW-D)/(TAW-RAW)$; after five days $D=139.8$ mm \ans $K_s=0.70$ \ans

\exbox{Crop example 8a: Cotton}{Cotton: a full bucket ($D=0$) drains by evapotranspiration at $K_cET_0=6.32$ mm\,d$^{-1}$ with no rain. After how many days does $K_s$ fall below 1, and what are $D$ and $K_s$ five days later?}\sol $TAW=187$, $RAW=122$ mm \ans $D$ reaches RAW after $19.2$ d \ans Then each day $D\leftarrow D+K_sK_cET_0$ with $K_s=(TAW-D)/(TAW-RAW)$; after five days $D=147.6$ mm \ans $K_s=0.60$ \ans

\exbox{Crop example 8a: Tomato}{Tomato: a full bucket ($D=0$) drains by evapotranspiration at $K_cET_0=6.32$ mm\,d$^{-1}$ with no rain. After how many days does $K_s$ fall below 1, and what are $D$ and $K_s$ five days later?}\sol $TAW=119$, $RAW=48$ mm \ans $D$ reaches RAW after $7.5$ d \ans Then each day $D\leftarrow D+K_sK_cET_0$ with $K_s=(TAW-D)/(TAW-RAW)$; after five days $D=74.1$ mm \ans $K_s=0.63$ \ans

\exbox{Crop example 8a: Potato}{Potato: a full bucket ($D=0$) drains by evapotranspiration at $K_cET_0=6.32$ mm\,d$^{-1}$ with no rain. After how many days does $K_s$ fall below 1, and what are $D$ and $K_s$ five days later?}\sol $TAW=85$, $RAW=30$ mm \ans $D$ reaches RAW after $4.7$ d \ans Then each day $D\leftarrow D+K_sK_cET_0$ with $K_s=(TAW-D)/(TAW-RAW)$; after five days $D=54.9$ mm \ans $K_s=0.54$ \ans

\exbox{Crop example 8a: Soybean}{Soybean: a full bucket ($D=0$) drains by evapotranspiration at $K_cET_0=6.32$ mm\,d$^{-1}$ with no rain. After how many days does $K_s$ fall below 1, and what are $D$ and $K_s$ five days later?}\sol $TAW=102$, $RAW=51$ mm \ans $D$ reaches RAW after $8.1$ d \ans Then each day $D\leftarrow D+K_sK_cET_0$ with $K_s=(TAW-D)/(TAW-RAW)$; after five days $D=75.7$ mm \ans $K_s=0.52$ \ans

\exbox{Crop example 8a: Sugar beet}{Sugar beet: a full bucket ($D=0$) drains by evapotranspiration at $K_cET_0=6.60$ mm\,d$^{-1}$ with no rain. After how many days does $K_s$ fall below 1, and what are $D$ and $K_s$ five days later?}\sol $TAW=153$, $RAW=84$ mm \ans $D$ reaches RAW after $12.8$ d \ans Then each day $D\leftarrow D+K_sK_cET_0$ with $K_s=(TAW-D)/(TAW-RAW)$; after five days $D=111.4$ mm \ans $K_s=0.60$ \ans

\exbox{Crop example 8a: Sunflower}{Sunflower: a full bucket ($D=0$) drains by evapotranspiration at $K_cET_0=6.05$ mm\,d$^{-1}$ with no rain. After how many days does $K_s$ fall below 1, and what are $D$ and $K_s$ five days later?}\sol $TAW=136$, $RAW=61$ mm \ans $D$ reaches RAW after $10.1$ d \ans Then each day $D\leftarrow D+K_sK_cET_0$ with $K_s=(TAW-D)/(TAW-RAW)$; after five days $D=86.9$ mm \ans $K_s=0.66$ \ans

\exbox{Crop example 8a: Onion}{Onion: a full bucket ($D=0$) drains by evapotranspiration at $K_cET_0=5.78$ mm\,d$^{-1}$ with no rain. After how many days does $K_s$ fall below 1, and what are $D$ and $K_s$ five days later?}\sol $TAW=60$, $RAW=18$ mm \ans $D$ reaches RAW after $3.1$ d \ans Then each day $D\leftarrow D+K_sK_cET_0$ with $K_s=(TAW-D)/(TAW-RAW)$; after five days $D=39.8$ mm \ans $K_s=0.47$ \ans

\exbox{Crop example 8a: Alfalfa}{Alfalfa: a full bucket ($D=0$) drains by evapotranspiration at $K_cET_0=6.32$ mm\,d$^{-1}$ with no rain. After how many days does $K_s$ fall below 1, and what are $D$ and $K_s$ five days later?}\sol $TAW=170$, $RAW=94$ mm \ans $D$ reaches RAW after $14.8$ d \ans Then each day $D\leftarrow D+K_sK_cET_0$ with $K_s=(TAW-D)/(TAW-RAW)$; after five days $D=120.3$ mm \ans $K_s=0.65$ \ans

}{}
\IfFileExists{extra/teny08.tex}{\subsection{Worked solutions on the ten soils (third scenario)}
The same questions are repeated for different working conditions, to show how the answer moves with them.

\exbox{Ten-soil example 8c: Coarse loam (S01)}{Coarse loam (hydrologic group B, $CN=75$): a storm of 60 mm falls on the soil. Find the runoff, the effective rainfall and the share that fills the bucket of $TAW=90$ mm if the profile was at 50 per cent depletion.}\sol $S=25\,400/75-254=84.7$ mm; $I_a=0.2S=16.9$ mm \ans $Q=(60-16.9)^2/(60-16.9+84.7)=14.5$ mm \ans Effective rain $60-14.5=45.5$ mm \ans The bucket holds $RAW=45$ mm of room, so it overfills: drainage 0.5 mm \ans

\exbox{Ten-soil example 8c: Coarse sand (S02)}{Coarse sand (hydrologic group A, $CN=65$): a storm of 60 mm falls on the soil. Find the runoff, the effective rainfall and the share that fills the bucket of $TAW=30$ mm if the profile was at 50 per cent depletion.}\sol $S=25\,400/65-254=136.8$ mm; $I_a=0.2S=27.4$ mm \ans $Q=(60-27.4)^2/(60-27.4+136.8)=6.3$ mm \ans Effective rain $60-6.3=53.7$ mm \ans The bucket holds $RAW=15$ mm of room, so it overfills: drainage 38.7 mm \ans

\exbox{Ten-soil example 8c: Sand (S03)}{Sand (hydrologic group A, $CN=65$): a storm of 60 mm falls on the soil. Find the runoff, the effective rainfall and the share that fills the bucket of $TAW=36$ mm if the profile was at 50 per cent depletion.}\sol $S=25\,400/65-254=136.8$ mm; $I_a=0.2S=27.4$ mm \ans $Q=(60-27.4)^2/(60-27.4+136.8)=6.3$ mm \ans Effective rain $60-6.3=53.7$ mm \ans The bucket holds $RAW=18$ mm of room, so it overfills: drainage 35.7 mm \ans

\exbox{Ten-soil example 8c: Loamy sand (S04)}{Loamy sand (hydrologic group A, $CN=65$): a storm of 60 mm falls on the soil. Find the runoff, the effective rainfall and the share that fills the bucket of $TAW=54$ mm if the profile was at 50 per cent depletion.}\sol $S=25\,400/65-254=136.8$ mm; $I_a=0.2S=27.4$ mm \ans $Q=(60-27.4)^2/(60-27.4+136.8)=6.3$ mm \ans Effective rain $60-6.3=53.7$ mm \ans The bucket holds $RAW=27$ mm of room, so it overfills: drainage 26.7 mm \ans

\exbox{Ten-soil example 8c: Sandy loam (S05)}{Sandy loam (hydrologic group B, $CN=75$): a storm of 60 mm falls on the soil. Find the runoff, the effective rainfall and the share that fills the bucket of $TAW=72$ mm if the profile was at 50 per cent depletion.}\sol $S=25\,400/75-254=84.7$ mm; $I_a=0.2S=16.9$ mm \ans $Q=(60-16.9)^2/(60-16.9+84.7)=14.5$ mm \ans Effective rain $60-14.5=45.5$ mm \ans The bucket holds $RAW=36$ mm of room, so it overfills: drainage 9.5 mm \ans

\exbox{Ten-soil example 8c: Loam (S06)}{Loam (hydrologic group B, $CN=75$): a storm of 60 mm falls on the soil. Find the runoff, the effective rainfall and the share that fills the bucket of $TAW=102$ mm if the profile was at 50 per cent depletion.}\sol $S=25\,400/75-254=84.7$ mm; $I_a=0.2S=16.9$ mm \ans $Q=(60-16.9)^2/(60-16.9+84.7)=14.5$ mm \ans Effective rain $60-14.5=45.5$ mm \ans The bucket holds $RAW=51$ mm of room, so it absorbs all: drainage 0.0 mm \ans

\exbox{Ten-soil example 8c: Silt loam (S07)}{Silt loam (hydrologic group C, $CN=82$): a storm of 60 mm falls on the soil. Find the runoff, the effective rainfall and the share that fills the bucket of $TAW=120$ mm if the profile was at 50 per cent depletion.}\sol $S=25\,400/82-254=55.8$ mm; $I_a=0.2S=11.2$ mm \ans $Q=(60-11.2)^2/(60-11.2+55.8)=22.8$ mm \ans Effective rain $60-22.8=37.2$ mm \ans The bucket holds $RAW=60$ mm of room, so it absorbs all: drainage 0.0 mm \ans

\exbox{Ten-soil example 8c: Sandy clay loam (S08)}{Sandy clay loam (hydrologic group C, $CN=82$): a storm of 60 mm falls on the soil. Find the runoff, the effective rainfall and the share that fills the bucket of $TAW=84$ mm if the profile was at 50 per cent depletion.}\sol $S=25\,400/82-254=55.8$ mm; $I_a=0.2S=11.2$ mm \ans $Q=(60-11.2)^2/(60-11.2+55.8)=22.8$ mm \ans Effective rain $60-22.8=37.2$ mm \ans The bucket holds $RAW=42$ mm of room, so it absorbs all: drainage 0.0 mm \ans

\exbox{Ten-soil example 8c: Clay loam (S09)}{Clay loam (hydrologic group D, $CN=87$): a storm of 60 mm falls on the soil. Find the runoff, the effective rainfall and the share that fills the bucket of $TAW=96$ mm if the profile was at 50 per cent depletion.}\sol $S=25\,400/87-254=38.0$ mm; $I_a=0.2S=7.6$ mm \ans $Q=(60-7.6)^2/(60-7.6+38.0)=30.4$ mm \ans Effective rain $60-30.4=29.6$ mm \ans The bucket holds $RAW=48$ mm of room, so it absorbs all: drainage 0.0 mm \ans

\exbox{Ten-soil example 8c: Silty clay loam (S10)}{Silty clay loam (hydrologic group D, $CN=87$): a storm of 60 mm falls on the soil. Find the runoff, the effective rainfall and the share that fills the bucket of $TAW=102$ mm if the profile was at 50 per cent depletion.}\sol $S=25\,400/87-254=38.0$ mm; $I_a=0.2S=7.6$ mm \ans $Q=(60-7.6)^2/(60-7.6+38.0)=30.4$ mm \ans Effective rain $60-30.4=29.6$ mm \ans The bucket holds $RAW=51$ mm of room, so it absorbs all: drainage 0.0 mm \ans

}{}
\IfFileExists{extra/tenby08.tex}{\subsection{Worked solutions on ten crops (hot climate)}
The ten crops of the library differ in their demand, their roots and their sensitivity. The same type of question is answered for each of them at $ET_0=7.5$ mm\,d$^{-1}$; the data are typical values of the crops (Doorenbos--Kassam and FAO 56).

\exbox{Crop example 8c: Maize}{Maize: a full bucket ($D=0$) drains by evapotranspiration at $K_cET_0=9.00$ mm\,d$^{-1}$ with no rain. After how many days does $K_s$ fall below 1, and what are $D$ and $K_s$ five days later?}\sol $TAW=170$, $RAW=94$ mm \ans $D$ reaches RAW after $10.4$ d \ans Then each day $D\leftarrow D+K_sK_cET_0$ with $K_s=(TAW-D)/(TAW-RAW)$; after five days $D=129.1$ mm \ans $K_s=0.53$ \ans

\exbox{Crop example 8c: Wheat}{Wheat: a full bucket ($D=0$) drains by evapotranspiration at $K_cET_0=8.62$ mm\,d$^{-1}$ with no rain. After how many days does $K_s$ fall below 1, and what are $D$ and $K_s$ five days later?}\sol $TAW=204$, $RAW=112$ mm \ans $D$ reaches RAW after $13.0$ d \ans Then each day $D\leftarrow D+K_sK_cET_0$ with $K_s=(TAW-D)/(TAW-RAW)$; after five days $D=147.9$ mm \ans $K_s=0.61$ \ans

\exbox{Crop example 8c: Cotton}{Cotton: a full bucket ($D=0$) drains by evapotranspiration at $K_cET_0=8.62$ mm\,d$^{-1}$ with no rain. After how many days does $K_s$ fall below 1, and what are $D$ and $K_s$ five days later?}\sol $TAW=187$, $RAW=122$ mm \ans $D$ reaches RAW after $14.1$ d \ans Then each day $D\leftarrow D+K_sK_cET_0$ with $K_s=(TAW-D)/(TAW-RAW)$; after five days $D=154.7$ mm \ans $K_s=0.49$ \ans

\exbox{Crop example 8c: Tomato}{Tomato: a full bucket ($D=0$) drains by evapotranspiration at $K_cET_0=8.62$ mm\,d$^{-1}$ with no rain. After how many days does $K_s$ fall below 1, and what are $D$ and $K_s$ five days later?}\sol $TAW=119$, $RAW=48$ mm \ans $D$ reaches RAW after $5.5$ d \ans Then each day $D\leftarrow D+K_sK_cET_0$ with $K_s=(TAW-D)/(TAW-RAW)$; after five days $D=81.5$ mm \ans $K_s=0.53$ \ans

\exbox{Crop example 8c: Potato}{Potato: a full bucket ($D=0$) drains by evapotranspiration at $K_cET_0=8.62$ mm\,d$^{-1}$ with no rain. After how many days does $K_s$ fall below 1, and what are $D$ and $K_s$ five days later?}\sol $TAW=85$, $RAW=30$ mm \ans $D$ reaches RAW after $3.4$ d \ans Then each day $D\leftarrow D+K_sK_cET_0$ with $K_s=(TAW-D)/(TAW-RAW)$; after five days $D=61.4$ mm \ans $K_s=0.43$ \ans

\exbox{Crop example 8c: Soybean}{Soybean: a full bucket ($D=0$) drains by evapotranspiration at $K_cET_0=8.62$ mm\,d$^{-1}$ with no rain. After how many days does $K_s$ fall below 1, and what are $D$ and $K_s$ five days later?}\sol $TAW=102$, $RAW=51$ mm \ans $D$ reaches RAW after $5.9$ d \ans Then each day $D\leftarrow D+K_sK_cET_0$ with $K_s=(TAW-D)/(TAW-RAW)$; after five days $D=81.8$ mm \ans $K_s=0.40$ \ans

\exbox{Crop example 8c: Sugar beet}{Sugar beet: a full bucket ($D=0$) drains by evapotranspiration at $K_cET_0=9.00$ mm\,d$^{-1}$ with no rain. After how many days does $K_s$ fall below 1, and what are $D$ and $K_s$ five days later?}\sol $TAW=153$, $RAW=84$ mm \ans $D$ reaches RAW after $9.4$ d \ans Then each day $D\leftarrow D+K_sK_cET_0$ with $K_s=(TAW-D)/(TAW-RAW)$; after five days $D=118.8$ mm \ans $K_s=0.50$ \ans

\exbox{Crop example 8c: Sunflower}{Sunflower: a full bucket ($D=0$) drains by evapotranspiration at $K_cET_0=8.25$ mm\,d$^{-1}$ with no rain. After how many days does $K_s$ fall below 1, and what are $D$ and $K_s$ five days later?}\sol $TAW=136$, $RAW=61$ mm \ans $D$ reaches RAW after $7.4$ d \ans Then each day $D\leftarrow D+K_sK_cET_0$ with $K_s=(TAW-D)/(TAW-RAW)$; after five days $D=94.3$ mm \ans $K_s=0.56$ \ans

\exbox{Crop example 8c: Onion}{Onion: a full bucket ($D=0$) drains by evapotranspiration at $K_cET_0=7.88$ mm\,d$^{-1}$ with no rain. After how many days does $K_s$ fall below 1, and what are $D$ and $K_s$ five days later?}\sol $TAW=60$, $RAW=18$ mm \ans $D$ reaches RAW after $2.3$ d \ans Then each day $D\leftarrow D+K_sK_cET_0$ with $K_s=(TAW-D)/(TAW-RAW)$; after five days $D=44.9$ mm \ans $K_s=0.35$ \ans

\exbox{Crop example 8c: Alfalfa}{Alfalfa: a full bucket ($D=0$) drains by evapotranspiration at $K_cET_0=8.62$ mm\,d$^{-1}$ with no rain. After how many days does $K_s$ fall below 1, and what are $D$ and $K_s$ five days later?}\sol $TAW=170$, $RAW=94$ mm \ans $D$ reaches RAW after $10.8$ d \ans Then each day $D\leftarrow D+K_sK_cET_0$ with $K_s=(TAW-D)/(TAW-RAW)$; after five days $D=127.9$ mm \ans $K_s=0.55$ \ans

}{}
\IfFileExists{extra/tenx08.tex}{\section{Worked solutions on the ten soils (second scenario)}
The same questions are repeated for different working conditions, to show how the answer moves with them.

\exbox{Ten-soil example 8b: Coarse loam (S01)}{Coarse loam (hydrologic group B, $CN=75$): a storm of 60 mm falls on the soil. Find the runoff, the effective rainfall and the share that fills the bucket of $TAW=90$ mm if the profile was at 50 per cent depletion.}\sol $S=25\,400/75-254=84.7$ mm; $I_a=0.2S=16.9$ mm \ans $Q=(60-16.9)^2/(60-16.9+84.7)=14.5$ mm \ans Effective rain $60-14.5=45.5$ mm \ans The bucket holds $RAW=45$ mm of room, so it overfills: drainage 0.5 mm \ans

\exbox{Ten-soil example 8b: Coarse sand (S02)}{Coarse sand (hydrologic group A, $CN=65$): a storm of 60 mm falls on the soil. Find the runoff, the effective rainfall and the share that fills the bucket of $TAW=30$ mm if the profile was at 50 per cent depletion.}\sol $S=25\,400/65-254=136.8$ mm; $I_a=0.2S=27.4$ mm \ans $Q=(60-27.4)^2/(60-27.4+136.8)=6.3$ mm \ans Effective rain $60-6.3=53.7$ mm \ans The bucket holds $RAW=15$ mm of room, so it overfills: drainage 38.7 mm \ans

\exbox{Ten-soil example 8b: Sand (S03)}{Sand (hydrologic group A, $CN=65$): a storm of 60 mm falls on the soil. Find the runoff, the effective rainfall and the share that fills the bucket of $TAW=36$ mm if the profile was at 50 per cent depletion.}\sol $S=25\,400/65-254=136.8$ mm; $I_a=0.2S=27.4$ mm \ans $Q=(60-27.4)^2/(60-27.4+136.8)=6.3$ mm \ans Effective rain $60-6.3=53.7$ mm \ans The bucket holds $RAW=18$ mm of room, so it overfills: drainage 35.7 mm \ans

\exbox{Ten-soil example 8b: Loamy sand (S04)}{Loamy sand (hydrologic group A, $CN=65$): a storm of 60 mm falls on the soil. Find the runoff, the effective rainfall and the share that fills the bucket of $TAW=54$ mm if the profile was at 50 per cent depletion.}\sol $S=25\,400/65-254=136.8$ mm; $I_a=0.2S=27.4$ mm \ans $Q=(60-27.4)^2/(60-27.4+136.8)=6.3$ mm \ans Effective rain $60-6.3=53.7$ mm \ans The bucket holds $RAW=27$ mm of room, so it overfills: drainage 26.7 mm \ans

\exbox{Ten-soil example 8b: Sandy loam (S05)}{Sandy loam (hydrologic group B, $CN=75$): a storm of 60 mm falls on the soil. Find the runoff, the effective rainfall and the share that fills the bucket of $TAW=72$ mm if the profile was at 50 per cent depletion.}\sol $S=25\,400/75-254=84.7$ mm; $I_a=0.2S=16.9$ mm \ans $Q=(60-16.9)^2/(60-16.9+84.7)=14.5$ mm \ans Effective rain $60-14.5=45.5$ mm \ans The bucket holds $RAW=36$ mm of room, so it overfills: drainage 9.5 mm \ans

\exbox{Ten-soil example 8b: Loam (S06)}{Loam (hydrologic group B, $CN=75$): a storm of 60 mm falls on the soil. Find the runoff, the effective rainfall and the share that fills the bucket of $TAW=102$ mm if the profile was at 50 per cent depletion.}\sol $S=25\,400/75-254=84.7$ mm; $I_a=0.2S=16.9$ mm \ans $Q=(60-16.9)^2/(60-16.9+84.7)=14.5$ mm \ans Effective rain $60-14.5=45.5$ mm \ans The bucket holds $RAW=51$ mm of room, so it absorbs all: drainage 0.0 mm \ans

\exbox{Ten-soil example 8b: Silt loam (S07)}{Silt loam (hydrologic group C, $CN=82$): a storm of 60 mm falls on the soil. Find the runoff, the effective rainfall and the share that fills the bucket of $TAW=120$ mm if the profile was at 50 per cent depletion.}\sol $S=25\,400/82-254=55.8$ mm; $I_a=0.2S=11.2$ mm \ans $Q=(60-11.2)^2/(60-11.2+55.8)=22.8$ mm \ans Effective rain $60-22.8=37.2$ mm \ans The bucket holds $RAW=60$ mm of room, so it absorbs all: drainage 0.0 mm \ans

\exbox{Ten-soil example 8b: Sandy clay loam (S08)}{Sandy clay loam (hydrologic group C, $CN=82$): a storm of 60 mm falls on the soil. Find the runoff, the effective rainfall and the share that fills the bucket of $TAW=84$ mm if the profile was at 50 per cent depletion.}\sol $S=25\,400/82-254=55.8$ mm; $I_a=0.2S=11.2$ mm \ans $Q=(60-11.2)^2/(60-11.2+55.8)=22.8$ mm \ans Effective rain $60-22.8=37.2$ mm \ans The bucket holds $RAW=42$ mm of room, so it absorbs all: drainage 0.0 mm \ans

\exbox{Ten-soil example 8b: Clay loam (S09)}{Clay loam (hydrologic group D, $CN=87$): a storm of 60 mm falls on the soil. Find the runoff, the effective rainfall and the share that fills the bucket of $TAW=96$ mm if the profile was at 50 per cent depletion.}\sol $S=25\,400/87-254=38.0$ mm; $I_a=0.2S=7.6$ mm \ans $Q=(60-7.6)^2/(60-7.6+38.0)=30.4$ mm \ans Effective rain $60-30.4=29.6$ mm \ans The bucket holds $RAW=48$ mm of room, so it absorbs all: drainage 0.0 mm \ans

\exbox{Ten-soil example 8b: Silty clay loam (S10)}{Silty clay loam (hydrologic group D, $CN=87$): a storm of 60 mm falls on the soil. Find the runoff, the effective rainfall and the share that fills the bucket of $TAW=102$ mm if the profile was at 50 per cent depletion.}\sol $S=25\,400/87-254=38.0$ mm; $I_a=0.2S=7.6$ mm \ans $Q=(60-7.6)^2/(60-7.6+38.0)=30.4$ mm \ans Effective rain $60-30.4=29.6$ mm \ans The bucket holds $RAW=51$ mm of room, so it absorbs all: drainage 0.0 mm \ans

}{}
\IfFileExists{extra/tenbx08.tex}{\subsection{Worked solutions on ten crops (cool climate)}
The ten crops of the library differ in their demand, their roots and their sensitivity. The same type of question is answered for each of them at $ET_0=4.0$ mm\,d$^{-1}$; the data are typical values of the crops (Doorenbos--Kassam and FAO 56).

\exbox{Crop example 8b: Maize}{Maize: a full bucket ($D=0$) drains by evapotranspiration at $K_cET_0=4.80$ mm\,d$^{-1}$ with no rain. After how many days does $K_s$ fall below 1, and what are $D$ and $K_s$ five days later?}\sol $TAW=170$, $RAW=94$ mm \ans $D$ reaches RAW after $19.5$ d \ans Then each day $D\leftarrow D+K_sK_cET_0$ with $K_s=(TAW-D)/(TAW-RAW)$; after five days $D=114.7$ mm \ans $K_s=0.72$ \ans

\exbox{Crop example 8b: Wheat}{Wheat: a full bucket ($D=0$) drains by evapotranspiration at $K_cET_0=4.60$ mm\,d$^{-1}$ with no rain. After how many days does $K_s$ fall below 1, and what are $D$ and $K_s$ five days later?}\sol $TAW=204$, $RAW=112$ mm \ans $D$ reaches RAW after $24.4$ d \ans Then each day $D\leftarrow D+K_sK_cET_0$ with $K_s=(TAW-D)/(TAW-RAW)$; after five days $D=133.0$ mm \ans $K_s=0.77$ \ans

\exbox{Crop example 8b: Cotton}{Cotton: a full bucket ($D=0$) drains by evapotranspiration at $K_cET_0=4.60$ mm\,d$^{-1}$ with no rain. After how many days does $K_s$ fall below 1, and what are $D$ and $K_s$ five days later?}\sol $TAW=187$, $RAW=122$ mm \ans $D$ reaches RAW after $26.4$ d \ans Then each day $D\leftarrow D+K_sK_cET_0$ with $K_s=(TAW-D)/(TAW-RAW)$; after five days $D=141.5$ mm \ans $K_s=0.69$ \ans

\exbox{Crop example 8b: Tomato}{Tomato: a full bucket ($D=0$) drains by evapotranspiration at $K_cET_0=4.60$ mm\,d$^{-1}$ with no rain. After how many days does $K_s$ fall below 1, and what are $D$ and $K_s$ five days later?}\sol $TAW=119$, $RAW=48$ mm \ans $D$ reaches RAW after $10.3$ d \ans Then each day $D\leftarrow D+K_sK_cET_0$ with $K_s=(TAW-D)/(TAW-RAW)$; after five days $D=67.8$ mm \ans $K_s=0.72$ \ans

\exbox{Crop example 8b: Potato}{Potato: a full bucket ($D=0$) drains by evapotranspiration at $K_cET_0=4.60$ mm\,d$^{-1}$ with no rain. After how many days does $K_s$ fall below 1, and what are $D$ and $K_s$ five days later?}\sol $TAW=85$, $RAW=30$ mm \ans $D$ reaches RAW after $6.5$ d \ans Then each day $D\leftarrow D+K_sK_cET_0$ with $K_s=(TAW-D)/(TAW-RAW)$; after five days $D=49.2$ mm \ans $K_s=0.65$ \ans

\exbox{Crop example 8b: Soybean}{Soybean: a full bucket ($D=0$) drains by evapotranspiration at $K_cET_0=4.60$ mm\,d$^{-1}$ with no rain. After how many days does $K_s$ fall below 1, and what are $D$ and $K_s$ five days later?}\sol $TAW=102$, $RAW=51$ mm \ans $D$ reaches RAW after $11.1$ d \ans Then each day $D\leftarrow D+K_sK_cET_0$ with $K_s=(TAW-D)/(TAW-RAW)$; after five days $D=70.2$ mm \ans $K_s=0.62$ \ans

\exbox{Crop example 8b: Sugar beet}{Sugar beet: a full bucket ($D=0$) drains by evapotranspiration at $K_cET_0=4.80$ mm\,d$^{-1}$ with no rain. After how many days does $K_s$ fall below 1, and what are $D$ and $K_s$ five days later?}\sol $TAW=153$, $RAW=84$ mm \ans $D$ reaches RAW after $17.5$ d \ans Then each day $D\leftarrow D+K_sK_cET_0$ with $K_s=(TAW-D)/(TAW-RAW)$; after five days $D=105.0$ mm \ans $K_s=0.70$ \ans

\exbox{Crop example 8b: Sunflower}{Sunflower: a full bucket ($D=0$) drains by evapotranspiration at $K_cET_0=4.40$ mm\,d$^{-1}$ with no rain. After how many days does $K_s$ fall below 1, and what are $D$ and $K_s$ five days later?}\sol $TAW=136$, $RAW=61$ mm \ans $D$ reaches RAW after $13.9$ d \ans Then each day $D\leftarrow D+K_sK_cET_0$ with $K_s=(TAW-D)/(TAW-RAW)$; after five days $D=80.8$ mm \ans $K_s=0.74$ \ans

\exbox{Crop example 8b: Onion}{Onion: a full bucket ($D=0$) drains by evapotranspiration at $K_cET_0=4.20$ mm\,d$^{-1}$ with no rain. After how many days does $K_s$ fall below 1, and what are $D$ and $K_s$ five days later?}\sol $TAW=60$, $RAW=18$ mm \ans $D$ reaches RAW after $4.2$ d \ans Then each day $D\leftarrow D+K_sK_cET_0$ with $K_s=(TAW-D)/(TAW-RAW)$; after five days $D=35.0$ mm \ans $K_s=0.59$ \ans

\exbox{Crop example 8b: Alfalfa}{Alfalfa: a full bucket ($D=0$) drains by evapotranspiration at $K_cET_0=4.60$ mm\,d$^{-1}$ with no rain. After how many days does $K_s$ fall below 1, and what are $D$ and $K_s$ five days later?}\sol $TAW=170$, $RAW=94$ mm \ans $D$ reaches RAW after $20.3$ d \ans Then each day $D\leftarrow D+K_sK_cET_0$ with $K_s=(TAW-D)/(TAW-RAW)$; after five days $D=113.9$ mm \ans $K_s=0.73$ \ans

}{}
\IfFileExists{extra/ex08.tex}{\backhead{EXERCISES}{khred}
\begin{exlist}
\item A storm of 45 mm falls on $CN=80$. Find the runoff.\exans{10.9 mm}
\item With $TAW=100$, $RAW=50$, $D_{i-1}=48$, $ET_0=6$, $K_c=1.15$, no rain: find $K_s$, $ET_a$ and $D_i$.\exans{1.00; 6.90 mm; 54.90 mm}
\item Same soil, $D_{i-1}=80$: find $K_s$, $ET_a$ and $D_i$.\exans{0.40; 2.76 mm; 82.76 mm}
\item $D_{i-1}=20$, rain 35 mm, $ET_a=5$: find $D_i$ and drainage.\exans{0; 10 mm}
\item A season gives $I=300$, $ET_a$ irrigated 540, $ET_a$ rainfed 430. Find IE.\exans{0.367}
\item Yield 9.4 Mg\,ha$^{-1}$, $ET_a=480$ mm: find WP.\exans{1.96 kg m$^{-3}$}
\item Find the curve number that gives $S=120$ mm.\exans{68}
\item A bucket with TAW 90 mm starts full (D = 0). ET = 5 mm/d, no rain. After how many days is $D=RAW=45$ mm ($p=0.5$)?\exans{9 d}
\item Relative yield for $ET_a/ET_c=0.85$, $K_y=1.2$.\exans{0.82}
\item State the assumptions of the tipping bucket in three lines.\exans{one layer; drainage the same day; no capillary rise}
\end{exlist}
}{}
\IfFileExists{extra/dyk08.tex}{\dykband{DO YOU KNOW ?}
\begin{exlist}
\item Rank soil-water models from the tipping bucket to Richards' equation and say when each is enough.
\item Write the daily update of the bucket with the stress coefficient.
\item What does the curve number represent, and why does it depend on antecedent moisture?
\item Why can a lower trigger give no extra yield?
\item What are the indicators of the dashboard?
\item Define the irrigation efficiency based on extra evapotranspiration.
\item Why is the bucket poor in a shallow soil with a water table?
\item How would you check that a bucket program closes the balance?
\end{exlist}
}{}
\IfFileExists{extra/obj08.tex}{\dykband{OBJECTIVE TYPE QUESTIONS}
\begin{exlist}
\item The tipping bucket drains when\opt{the soil is below wilting}{the store exceeds field capacity}{there is no rain}{at night}
\item A larger CN means\opt{less runoff}{more runoff}{no runoff}{more infiltration}
\item $K_s=1$ when\opt{$D>RAW$}{$D\le RAW$}{$D=TAW$}{always zero}
\item The initial abstraction is\opt{$0.2S$}{$S/2$}{$5S$}{$S^2$}
\item Water productivity is\opt{$Y/ET_a$}{$ET_a/Y$}{$P/Y$}{$I/Y$}
\item The rainfed twin is used to find\opt{the soil type}{the extra effect of irrigation}{the weather}{nothing}
\item A trigger above $p$ causes\opt{no stress}{stress}{more rain}{drainage only}
\item The cascading bucket has\opt{one store}{several stores in series}{no stores}{no drainage}
\item Engine 2 of the application is\opt{a Richards solver}{a daily bucket}{a regression}{a satellite}
\item The dashboard closes the balance with\opt{the storage change}{the price}{the wind}{the albedo}
\end{exlist}
\objkey{1.~(b)\quad 2.~(b)\quad 3.~(b)\quad 4.~(a)\quad 5.~(a)\quad 6.~(b)\quad 7.~(b)\quad 8.~(b)\quad 9.~(b)\quad 10.~(a)}
}{}
\swaflvideo{08}
